\DocumentMetadata{
  lang = en,
  pdfversion = 2.0,
  pdfstandard = ua-2,
  tagging = on,
  tagging-setup = {
    math/setup = {mathml-SE},
    table/header-rows = 1
  }
}

\documentclass[aspectratio=169]{ltx-talk}

\usepackage{amsmath}

\tagpdfsetup{
  role/new-tag=frametitle/H1,
  math/alt/use
}

\title{Econ 905 -- Lecture 1}
\author{Your Name}
\institute{Kansas State University}
\date{Fall 2025}

\hypersetup{
  pdftitle={Econ 905 -- Lecture 1},
  pdfauthor={Your Name},
}

\begin{document}

\begin{frame}
  \titlepage
\end{frame}

\begin{frame}
\frametitle{Introduction}
We first deal with the most simple kind of macroeconomic model consisting of
a representative firm and a representative consumer. All agents will
optimize. They maximize some objective function subject to the constraints
they face. We will take as fundamental:

\begin{itemize}
\item the preferences of consumers

\item the endowments of resources available to consumers

\item the technology available to firms

\item market structure.
\end{itemize}

Assuming \textit{optimizing behavior} and a \textit{defining an equilibrium}
allows us to use the model to make predictions.

\textit{Preferences}

There is one period (a \textit{static} model as opposed to a \textit{dynamic
model}) and there are $N$ consumers. These consumers have identical
preference given by $u(c,\ell )$ where $c$ is consumption and $\ell $ is
leisure. We set $N=1$ with no loss of generality. This is our representative
agent.
\end{frame}

\begin{frame}
\frametitle{The utility function}
The utility function, $u(c,\ell ),$ has the following features

\begin{itemize}
\item strictly increasing in each argument

\item strictly concave

\item twice differentiable

\item $lim_{c\rightarrow 0}\,u_{1}(c,\ell )=\infty ,$ $\ell >0$ and $%
lim_{\ell \rightarrow 0}\,u_{2}(c,\ell )=\infty ,$ $c>0.$
\end{itemize}

\textit{Endowments}

Each consumer is endowed with one unit of time to be allocated between work
and leisure. Also, each consumer owns $k_{0}$ units of capital (also an
endowment) which can be rented to firms.

\textit{Technology}

There are $M$ firms. Each can produce consumption goods according to 
\begin{equation*}
y=zf(k,n)
\end{equation*}%
where $y$ is output, $k$ is the capital input, $n$ is the labor input, $z$
is a scalar.
\end{frame}

\begin{frame}
\frametitle{The production function}
The production function, $z$ $f(k,n),$ has the following features

\begin{itemize}
\item increasing in both arguments

\item strictly quasiconcave
\par\vspace{0.4em}{\footnotesize\textit{Note:} \textit{A} function $f$ a quasiconcave
function iff for any pair of distinct points $u$ and $v$ in the domain of $f$%
, and for $0<\theta <1$%
\begin{equation*}
f\left( v\right) \geq f\left( u\right) \Rightarrow f\left[ \theta u+\left(
1-\theta \right) v\right] \geq f\left( u\right)
\end{equation*}%
(strict) concavity implies (strict) quasiconcavity, but not vice-versa. Thus
quasi-concavity is a generalization of the notion of concavity and a
somewhat weaker condition. Any monotonic, non-decreasing function of a
quasiconcave function---and therefore of a concave function---is
quasi-concave.}


\item twice differentiable

\item homogeneous of degree 1
\par\vspace{0.4em}{\footnotesize\textit{Note:} %
This implies constant returns to scale, so that $\lambda y=zf(\lambda
k,\lambda n).$}


\item $lim_{k\rightarrow 0}\,f_{1}(k,n)=\infty ,$ $lim_{k\rightarrow \infty
}\,f_{1}(k,n)=0,$ $lim_{n\rightarrow 0}\,f_{2}(k,n)=\infty ,$ $%
lim_{n\rightarrow \infty }\,f_{2}(k,n)=0.$
\par\vspace{0.4em}{\footnotesize\textit{Note:} %
The notation $u_{1}(c,\ell )$ means the derivative of $u(c,\ell )$ with
respect to its first argument. Other derivatives are similarly defined.}
.
\end{itemize}

We set $M=1$ with no loss of generality.
\end{frame}

\begin{frame}
\frametitle{Market Structure}
We will focus on a \textit{competitive equilibrium} where all agents are
price takers. We can only determine relative prices so the price of one good
can arbitrarily be set to 1. We call this good the \textit{numeraire}. Treat
the consumption good as the numeraire. We have 3 markets:

\begin{enumerate}
\item consumption goods

\item leisure (labor): the price of leisure in units of consumption is $%
\omega $

\item rental services of capital: the rental rate on capital is $r.$
\end{enumerate}
\end{frame}

\begin{frame}
\frametitle{Consumer{'}s Problem}
Each agent solves%
\begin{equation*}
\underset{c,\ell ,k_{s}}{\max }u(c,\ell )\text{ subject to}
\end{equation*}%
\begin{eqnarray*}
c &\leq &\omega \left( 1-\ell \right) +rk_{s} \\
0 &\leq &k_{s}\leq k_{o} \\
0 &\leq &\ell \leq 1 \\
c &\geq &0.
\end{eqnarray*}%
Here $k_{s}$ is the quantity of capital that the consumer rents to firm(s)
and since utility is increasing in consumption we must have $k_{s}=k_{0}$.
Our restrictions on the utility function assure that:

\begin{enumerate}
\item the nonnegativity constraints on consumption and leisure will not be
binding

\item in equilibrium we will have $0<\ell <1.$
\end{enumerate}
\end{frame}

\begin{frame}
\frametitle{Lagrangian }
\begin{equation*}
L=u(c,\ell )+\mu \left( rk_{o}+\omega \left( 1-\ell \right) -c\right)
\end{equation*}%
There is sufficient structure on the utility function to assure a unique
optimum characterized by the following focs:%
\begin{eqnarray*}
u_{1}(c,\ell ) &=&\mu \\
u_{2}(c,\ell ) &=&\mu \omega \\
rk_{o}+\omega \left( 1-\ell \right) &=&c.
\end{eqnarray*}%
The first two imply $\omega u_{1}(c,\ell )=u_{2}(c,\ell ).$ Substitute in
for $c$ from the third to find

\begin{equation*}
\omega u_{1}(rk_{o}+\omega \left( 1-\ell \right) ,\ell )=u_{2}(rk_{o}+\omega
\left( 1-\ell \right) ,\ell ).
\end{equation*}%
This solves (implicitly)\ for the desired quantity of leisure, $\ell $; in
terms of $r,k_{o},\omega $ (one equation and one unknown). It does not,
however, solve the model. We were given $k_{o}.$ We were not given $r$ and $%
\omega .$
\end{frame}

\begin{frame}
\frametitle{Firms{'} Problem}
Each firm chooses inputs of labor and capital
to maximize profits, treating $\omega $ and $r$ as fixed; i.e. each solves%
\begin{equation*}
\underset{k,n}{\max }\left[ zf\left( k,n\right) -\omega n-rk\right] .
\end{equation*}%
Profit maximization requires 
\begin{equation*}
\begin{array}{ll}
zf_{1}\left( k,n\right) )=r & zf_{2}\left( k,n\right) )=\omega .%
\end{array}%
\end{equation*}%
Since $f\left( k,n\right) )$ is homogeneous of degree one, Euler's law holds%

\par\vspace{0.4em}{\footnotesize\textit{Note:} \textit{Euler's law:} Suppose$f\left( \mathbf{x}\right) $ is
homogeneous of degree $g$, If we let $\mathbf{x}=\left[ x_{1},x_{2},...x_{M}%
\right] $ (i.e. a vector with M) elements, then $gf\left( \mathbf{x}\right) =%
\overset{M}{\underset{i=1}{\sum f_{xi}}}\left( \mathbf{x}\right) x_{i}.$}
.
For us this implies that 
\begin{equation*}
zf(k,n)=zf_{1}k+zf_{2}n.
\end{equation*}%
That is, profits are zero. This has two important consequences:

\begin{itemize}
\item We need not be concerned with who owns the firm or how profits are
distributed.

\item The optimal scale of the firm is indeterminate and the number of firms
is irrelevant so we set it equal to 1.
\end{itemize}
\end{frame}

\begin{frame}
\frametitle{Definition}
\ A competitive equilibrium in this economy is
a set of quantities $\left\{ c,y,\ell ,n,k,k_{s}\right\} $ and prices $%
\left\{ \omega ,r\right\} $ which satisfy the following properties:

\begin{enumerate}
\item Each consumer chooses $c,$ $k_{s}$ and $\ell $ optimally given $\omega 
$ and $r,$

\item The representative firm chooses $y$, $n$ and $k$ optimally given $%
\omega $ and $r,$

\item Markets clear.
\end{enumerate}

We have three markets. In our competitive equilibrium, market clearing
requires the following%
\begin{equation*}
\begin{array}{lll}
(1-\ell )=n, & y=c, & k_{0}=k%
\end{array}%
.
\end{equation*}%
Note: the total value of \textit{excess demand} across markets is%
\begin{equation*}
c-y+\omega \lbrack n-(1-\ell )]+r(k-k_{0}).
\end{equation*}
\end{frame}

\begin{frame}
\frametitle{Walras{'} law}
\begin{itemize}
\item \textit{Walras' law} states that the value of excess demand across
markets is always zero, and this then implies that, if there are $X$ markets
and $X-1$ of those markets are in equilibrium, then the additional market is
also in equilibrium.

\item Note above that if any two of the market clearing conditions are met
and excess demand is zero, the third is necessarily met so this is an
example of Walras' law\textit{.}

\item Thus we can drop one market-clearing condition in determining
competitive equilibrium prices and quantities. Lets eliminate $y=c$.
\end{itemize}
\end{frame}

\begin{frame}
\frametitle{The competitive equilibrium reduces to a system of five equations in the}
The competitive equilibrium reduces to a system of five equations in the
five unknowns: $\omega ,r,\ell ,k,n.$ These are

\begin{equation*}
\omega u_{1}(rk_{o}+\omega \left( 1-\ell \right) ,\ell )=u_{2}(rk_{o}+\omega
\left( 1-\ell \right) ,\ell )
\end{equation*}%
\begin{equation*}
\begin{array}{llll}
zf_{1}\left( k,n\right) )=r, & zf_{2}\left( k,n\right) )=\omega , & (1-\ell
)=n, & k_{0}=k.%
\end{array}%
\end{equation*}%
We can plug the 2nd through 5th of these equations into the first to get: 
\begin{equation*}
zf_{2}\left( k_{o},1-\ell \right) u_{1}(zf\left( k_{o},1-\ell \right) ,\ell
)=u_{2}(zf\left( k_{o},\left( 1-\ell \right) \right) ,\ell )
\end{equation*}%
or more succinctly%
\begin{equation*}
zf_{2}u_{1}=u_{2}
\end{equation*}%
which solves implicitly for equilibrium $\ell $. Knowing $\ell $ we
substitute in the following equations to obtain solutions for $\omega ,r,k,n$
and $c$ using%
\begin{equation*}
zf_{1}=r;~zf_{2}=\omega ,~n=(1-\ell ),~k_{0}=k,~c=zf\left( k_{0},(1-\ell
)\right) .
\end{equation*}
\end{frame}

\begin{frame}
\frametitle{A Pareto Optimum is some allocation such that there is no other}
A \textit{Pareto Optimum}\ is some allocation such that there is no other
allocation which some agents strictly prefer and which does not make any
agents worse off.
\par\vspace{0.4em}{\footnotesize\textit{Note:} %
An allocation is a production plan and a distribution of goods across
economic agents.}


Here, we have a single agent and need not worry about the allocation of
goods across agents. It helps to think in terms of a \textit{social planner}%
. A social planner is an 'agent' (theoretical construct outside the model)\
who can:

\begin{itemize}
\item `dictate' inputs to production by the representative firm

\item `force' the consumer to supply the appropriate quantity of labor

\item `distribute' consumption goods to the consumer.
\end{itemize}

The planner does this so as to make the consumer as well off as possible.
\end{frame}

\begin{frame}
\frametitle{The social planner determines a Pareto optimum by solving the}
The \textit{social planner determines a Pareto optimum }by solving the
following problem%
\begin{equation*}
\underset{c,\ell }{\max }u(c,\ell )\text{ }
\end{equation*}%
subject to%
\begin{equation*}
c=zf\left( k_{o},1-\ell \right) .
\end{equation*}%
That is, the social planner tries to maximize welfare given a \textit{%
resource constraint. }We can write this as 
\begin{equation*}
\underset{\ell }{\max }u(zf\left( k_{o},1-\ell \right) ,\ell )\text{.}
\end{equation*}%
Setting the first order condition (foc) equal to zero gives:%
\begin{equation*}
zf_{2}\left( k_{o},1-\ell \right) u_{1}(zf\left( k_{o},1-\ell \right) ,\ell
)=u_{2}(zf\left( k_{o},1-\ell \right) ,\ell ).
\end{equation*}%
\end{frame}

\begin{frame}
\frametitle{This the same result we got (in terms of ) in solving for a}
This the same result we got (in terms of $\ell $) in solving for a
competitive equilibrium. That is, \textit{the competitive equilibrium and
the Pareto optimum are identical here}. Further, since $u(c,\ell )$ is
strictly concave, and $f\left( k_{o},1-\ell \right) ,\ell )$ is strictly
quasiconcave there is a unique Pareto optimum, and the competitive
equilibrium is also unique.

\noindent In more general settings, it is true \textit{under some
restrictions} that the following hold:

\begin{enumerate}
\item A competitive equilibrium (with perfect competition) is Pareto optimal
(\textit{First Welfare Theorem}).

\item Any Pareto optimum can be supported as a competitive equilibrium (with
perfect competition) with an appropriate choice of endowments. (\textit{%
Second Welfare Theorem}).
\end{enumerate}

Speaking loosely, the assumptions required for (1) and (2) to hold include:%

\par\vspace{0.4em}{\footnotesize\textit{Note:} %
See Varian's text \textit{Microeconomic Analysis }for proofs and a more
concise statement.}


\begin{itemize}
\item the absence of externalities

\item completeness of markets (discussed in next set of notes)

\item absence of distorting taxes and expenditures (e.g. income taxes and
sales taxes).
\end{itemize}
\end{frame}

\begin{frame}
\frametitle{Some uses of solving the social planner{'}s problem}
Some uses of solving the social planner's problem:

\noindent \textbf{When it holds}, the equivalence of competitive equilibria
and Pareto optima in representative agent models can be very useful...we can
solve for a competitive equilibrium \textit{or} solve the social planner's
problem (which is often easier).

\noindent \textbf{When it does not hold}:

\begin{itemize}
\item We often can use the Pareto optimal allocation to get insights into
proper government policy.

\item We often can use the Pareto optimal allocation to get a measure of the
cost of inefficiency.
\end{itemize}
\end{frame}

\begin{frame}
\frametitle{An example}
Assign particular functional forms. Preferences are 
\begin{equation*}
u\left( c,\ell \right) =\frac{c^{1-\gamma }-1}{1-\gamma }+\frac{\ell
^{1-\gamma }-1}{1-\gamma }
\end{equation*}%
where $\gamma >0$ measures the degree of curvature.
\par\noindent\textbf{Figure (plot omitted).} $\gamma =.5$ lighter solid line;$\ \gamma =1$
dashed line; $\gamma =1.5$ darker solid line.\par


This is a \textquotedblleft constant relative risk
aversion\textquotedblright\ (CRRA) utility function. At $\gamma =1$ we
define $u\left( c,\ell \right) =\ln c+\ln \ell .$
\par\vspace{0.4em}{\footnotesize\textit{Note:} %
The rationale for this is that using L'Hospital's Rule $\underset{\gamma
\rightarrow 1.}{\lim }\frac{c^{1-\gamma }-1}{1-\gamma }=\ln c.$ See Chiang's
text \textit{Fundamental Methods of Mathematical Economics}.}
\end{frame}

\begin{frame}
\frametitle{Production Cobb-Douglas}
Production \textit{Cobb-Douglas}:%
\begin{equation*}
f\left( k,n\right) =zk^{\alpha }n^{1-\alpha }\text{ }0<\alpha <1
\end{equation*}%
The social planner's problem is easiest here. The social planner solves:%
\begin{equation*}
u\left( \ell \right) =\frac{\left( zk^{\alpha }\left( 1-\ell \right)
^{1-\alpha }\right) ^{1-\gamma }-1}{1-\gamma }+\frac{\ell ^{1-\gamma }-1}{%
1-\gamma }
\end{equation*}%
Setting the foc equal to zero gives:%
\begin{equation*}
\left[ zk_{o}^{\alpha }\left( 1-\ell \right) ^{1-\alpha }\right] ^{-\gamma
}\left( 1-\alpha \right) zk_{o}^{\alpha }\left( 1-\ell \right) ^{-\alpha
}=\ell ^{-\gamma }
\end{equation*}%
This is one equation and one unknown.
\end{frame}

\begin{frame}
\frametitle{We can easily solve for it numerically to see how the labor supply}
We can easily solve for it \textit{numerically} to see how the labor supply
depends on $z$. We do ths for three different valules of $\gamma .$


\par\noindent\textbf{Figure (plot omitted).} $\gamma =.5$ lighter solid
line;$\ \gamma =1$ dashed line; $\gamma =1.5$ darker solid
line.\par
\end{frame}

\begin{frame}
\frametitle{Slide 18}
Lets look at a special case we can solve \textit{analytically}. Let $\gamma
=1$ i.e., the $ln$ case.%
\begin{equation*}
\left[ zk_{o}^{\alpha }\left( 1-\ell \right) ^{1-\alpha }\right] ^{-1}\left(
1-\alpha \right) zk_{o}^{\alpha }\left( 1-\ell \right) ^{-\alpha }=\ell ^{-1}
\end{equation*}%
or 
\begin{equation*}
\ell =\frac{1}{2-\alpha }.
\end{equation*}%
Here $\ell $ is not a function of $zk_{o}^{\alpha }.$ From this we can find
all other endogenous items. To look at the effects of $z$, we need to look
at the more general case above or we can consider a somewhat different set
of preferences. Lets set preferences to 
\begin{equation*}
u\left( c,\ell \right) =\frac{c^{1-\gamma }-1}{1-\gamma }+\ell .
\end{equation*}%
Point to ponder:\ For leisure this this does not satisfy all the assumptions
we made above. Utility is linear in leisure. Can we be assured that we will
not have a \textit{corner solution}?
\end{frame}

\begin{frame}
\frametitle{The social planner solves}
The social planner solves%
\begin{equation*}
\underset{\ell }{\max }\frac{\left[ zk_{o}^{\alpha }\left( 1-\ell \right)
^{1-\alpha }\right] ^{1-\gamma }-1}{1-\gamma }+\ell .
\end{equation*}%
Setting the foc equal to zero gives%
\begin{equation*}
\left[ zk_{o}^{\alpha }\left( 1-\ell \right) ^{1-\alpha }\right] ^{-\gamma
}\left( 1-\alpha \right) zk_{o}^{\alpha }\left( 1-\ell \right) ^{-\alpha }=1
\end{equation*}%
which gives (show this on your own)%
\begin{equation*}
\ell =1-\left[ \left( zk_{o}^{\alpha }\right) ^{1-\gamma }\left( 1-\alpha
\right) \right] ^{\frac{1}{\alpha +\left( 1-\alpha \right) \gamma }}.
\end{equation*}%
\end{frame}

\begin{frame}
\frametitle{Slide 20}
Knowing $\ell $ we can find $c$ from $c=zk_{o}^{\alpha }\left(
1-\ell \right) ^{1-\alpha }$ which gives (after some algebra not shown)%
\begin{equation*}
c=\left[ zk_{o}^{\alpha }\left( 1-\alpha \right) ^{\left( 1-\alpha \right) }%
\right] ^{\frac{1}{\alpha +\left( 1-\alpha \right) \gamma }}
\end{equation*}%
and since $zf_{2}\left( k_{o},\left( 1-\ell \right) \right) )=\ \omega $
(from the decentralized problem) $\allowbreak $

Note: Since $\ell $ cannot be less than zero, it is more complete to write 
\begin{equation*}
\ell =\max \left\{ 0,1-\left[ \left( zk_{o}^{\alpha }\right) ^{1-\gamma
}\left( 1-\alpha \right) \right] ^{\frac{1}{\alpha +\left( 1-\alpha \right)
\gamma }}\right\}
\end{equation*}%
\begin{equation*}
c=\left\{ 
\begin{array}{ccc}
\left[ zk_{o}^{\alpha }\left( 1-\alpha \right) ^{\left( 1-\alpha \right) }%
\right] ^{\frac{1}{\alpha +\left( 1-\alpha \right) \gamma }} & \text{if} & 
1> \left[ \left( zk_{o}^{\alpha }\right) ^{1-\gamma }\left( 1-\alpha \right) %
\right] ^{\frac{1}{\alpha +\left( 1-\alpha \right) \gamma }} \\ 
c=zk_{o}^{\alpha } & \text{if} & 1\leq \left[ \left( zk_{o}^{\alpha }\right)
^{1-\gamma }\left( 1-\alpha \right) \right] ^{\frac{1}{\alpha +\left(
1-\alpha \right) \gamma }}%
\end{array}%
\right. .
\end{equation*}%
In general, be sure to watch for corner solutions.
\end{frame}

\begin{frame}
\frametitle{Comparative statics.}
Both $c$ and $\omega $ are increasing in both $%
z$ and $k_{0}$. However, the effects on the quantity of leisure (and
therefore on employment) are ambiguous.

\begin{itemize}
\item With $\gamma $ $<1$ an increase in $z$ or in $k_{0}$ will decrease
leisure.

\item With $\gamma $ $>1$ an increase in $z$ or in $k_{0}$ will increase
leisure.

\item Why? There are opposing \textit{income and substitution effects. }
\end{itemize}

An empirical regularity is that an increase in productivity increases labor
employed (decreases leisure). A productivity shock, which increases the real
wage and output must result in a decrease in leisure (increase in
employment) in order for employment to be procyclical in our model as in the
data. The substitution effect must be large (in absolute value) relative to
the income effect. $\gamma $ $<1$ is sufficient for this in our example.
\end{frame}

\begin{frame}
\frametitle{An economy with a government}
We now introduce government into our model to analyze some simple fiscal
policy issues. We make some other simplifications to keep it clean. The
government

\begin{itemize}
\item makes purchases of consumption goods.

\item finances purchases through lump-sum taxes on the representative agent.

\item destroys the goods it purchases.
\end{itemize}

Let $g$ be the quantity of government purchases (exogenous) and $\tau $ be
total taxes. The government operates under a period \textit{balanced budget
constraint} so that $g=\tau .$

Output can be produced according to $y=zn$.
\end{frame}

\begin{frame}
\frametitle{Slide 23}
Here, the consumer's optimization problem is%
\begin{equation*}
\underset{c,\ell }{\max }u\left( c,\ell \right)
\end{equation*}%
subject to%
\begin{equation*}
c=\omega (1-\ell )-\tau .
\end{equation*}%
The foc for this maximization problem gives:%
\begin{equation*}
\omega u_{1}=u_{2}.
\end{equation*}%
The representative firm's profit maximization problem is%
\begin{equation*}
\underset{n}{\max }\left( z-\omega \right) n.
\end{equation*}%
To maximize profits the firm chooses 
\begin{equation*}
n=\left\{ 
\begin{array}{ccc}
-\infty & \text{if} & z<\omega \\ 
\text{any quantity} & \text{if} & z=\omega \\ 
\infty & \text{if} & z>\omega .%
\end{array}%
\right.
\end{equation*}%
\end{frame}

\begin{frame}
\frametitle{Slide 24}
Now a \textit{competitive equilibrium} consists of five quantities $\left(
y,c,\ell ,n,\tau \right) $ and a price $\left( \omega \right) $ which
satisfy the following:

\begin{enumerate}
\item The representative consumer chooses $c$ and $\ell $ to maximize
utility, given $\omega $ and $\tau $.

\item The representative firm chooses $n$ to maximize profits, given $\omega 
$.

\item Markets for consumption goods and labor clear.

\item The government budget constraint is satisfied.
\end{enumerate}

Point to ponder: How do we know there is only 1 price to solve for?
\end{frame}

\begin{frame}
\frametitle{The competitive equilibrium and the Pareto optimum are equivalent...}
The competitive equilibrium and the Pareto optimum are equivalent here, (no
distortionary taxes or spending). The social planner's problem is 
\begin{equation*}
\underset{c,\ell }{\max }u\left( c,\ell \right)
\end{equation*}%
subject to 
\begin{equation*}
c+g=z(1-\ell ).
\end{equation*}

After substituting in for $c$, the foc is: 
\begin{equation*}
zu_{1}\left[ \left( z\left( 1-\ell \right) -g,\ell \right) \right]
=u_{2}\left( z\left( 1-\ell \right) -g,\ell \right) .
\end{equation*}%
Suppose preferences are given by 
\begin{equation*}
u\left( c,\ell \right) =\frac{c^{1-\gamma }-1}{1-\gamma }+\frac{\ell
^{1-\gamma }-1}{1-\gamma }.
\end{equation*}%
Substitute in $c=z(1-\ell )-g$ 
\begin{equation*}
u\left( \ell s\right) =\frac{\left[ z(1-\ell )-g\right] ^{1-\gamma }-1}{%
1-\gamma }+\frac{\ell ^{1-\gamma }-1}{1-\gamma }.
\end{equation*}%
\end{frame}

\begin{frame}
\frametitle{The foc is}
The foc is%
\begin{equation*}
z\left[ z(1-\ell )-g\right] ^{-\gamma }=\ell ^{-\gamma }
\end{equation*}%
which gives 
\begin{equation*}
\ell =\frac{z^{\frac{-1}{\gamma }}\left( z-g\right) }{1+z^{\frac{\gamma -1}{%
\gamma }}}
\end{equation*}%
and%
\begin{equation*}
c=z\left( 1-\frac{z^{\frac{-1}{\gamma }}\left( z-g\right) }{1+z^{\frac{%
\gamma -1}{\gamma }}}\right) -g
\end{equation*}%
or%
\begin{equation*}
c=\left( \frac{z+z^{\frac{\gamma -1}{\gamma }}g}{1+z^{\frac{\gamma -1}{%
\gamma }}}\right) -g.
\end{equation*}%
\end{frame}

\begin{frame}
\frametitle{Slide 27}
\begin{itemize}
\item Note $\frac{d\ell }{dg}<0$ and $\frac{dc}{dg}<0.$

\item Quantities of both goods will decrease.

\item There is \textit{crowding out} of private consumption.

\item The decrease in consumption is smaller than the increase in government
purchases, ($\frac{dc}{dg}>-1$) so that output increases.
\end{itemize}

Another lesson from this is that wasteful government spending does not make
the social planner's problem differ from the competitive equilibrium. (On
your own: solve for CE\ to verify). Why? As we have set up the problem, the
social planner takes the government spending as \textit{fundamental}%
.
\end{frame}

\begin{frame}
\frametitle{An economy with distortionary taxes}
}

Let us examine a representative agent economy where the consumer has a
utility function 
\begin{equation}
u=\ln c+\ln \ell  \label{u2}
\end{equation}%
where $c$ and $\ell $ are consumption and leisure. The consumer has an
endowment of 1 unit of time and $k_{0}$ units of capital. The representative
firm has a production technology given by $y=k^{\alpha }n^{1-\alpha }$ where 
$k$ and $n$ are capital and labor inputs. Suppose the government must have
expenditure equal to $g$ which is financed by taxing consumption at rate $%
\tau $. The social planner maximizes welfare subject to a \textit{resource
constraint}. The problem is 
\begin{equation*}
\underset{c,\ell }{\max }\ln c+\ln \ell
\end{equation*}%
subject to:%
\begin{equation}
c+g=k^{\alpha }\left( 1-\ell \right) ^{1-\alpha }  \notag
\end{equation}%
yielding 
\begin{equation*}
L=\ln c+\ln \ell +\lambda \left( k^{\alpha }\left( 1-\ell \right) ^{1-\alpha
}-c-g\right) .
\end{equation*}%
\end{frame}

\begin{frame}
\frametitle{The focs are}
The focs are%
\begin{eqnarray*}
\frac{1}{c} &=&\lambda \\
\frac{1}{\ell } &=&\lambda \left( 1-\alpha \right) k^{\alpha }\left( 1-\ell
\right) ^{-\alpha }
\end{eqnarray*}%
giving%
\begin{equation}
\frac{\frac{1}{c}}{\frac{1}{\ell }}=\frac{1}{\left( 1-\alpha \right)
k^{\alpha }\left( 1-\ell \right) ^{-\alpha }}.  \label{spp}
\end{equation}%
Using the resource constraint and simplifying gives 
\begin{equation}
\frac{\left( 1-\alpha \right) k^{\alpha }\left( 1-\ell \right) ^{-\alpha }}{%
\left( k^{\alpha }\left( 1-\ell \right) ^{1-\alpha }-g\right) }=\frac{1}{%
\ell }  \label{6}
\end{equation}%
Since the size of the capital stock is $k_{0}$ and $g$ is given, this can be
used to find $\ell .$
\end{frame}

\begin{frame}
\frametitle{Slide 30}
A competitive equilibrium is a set of quantities $\left\{ y,c,\ell ,n,k,\tau
\right\} $ prices $\left\{ \omega ,r\right\} $ such that

\begin{enumerate}
\item Each consumer chooses $c$ and $\ell $ to maximize (\ref{u2}) given $%
\omega ,r$, and $\tau $.

\item The representative firm chooses $k$ and $n$ to maximize profits given $%
\omega $ and $r.$

\item The government balances its budget.

\item Markets clear. That is:%
\begin{eqnarray*}
k &=&k_{0} \\
n+\ell &=&1 \\
c+g &=&y.
\end{eqnarray*}%
\end{enumerate}
\end{frame}

\begin{frame}
\frametitle{Solve for the competitive equilibrium and show that it is not Pareto}
Solve for the competitive equilibrium and show that it is not Pareto
Optimal.. The agent's problem is
\begin{equation*}
\underset{c,\ell }{\max }\ln c+\ln \ell
\end{equation*}%
subject to:%
\begin{equation*}
c\left( 1+\tau \right) \leq \left( \omega \left( 1-\ell \right)
+rk_{0}\right) .
\end{equation*}%
This gives%
\begin{equation*}
L=\ln c+\ln \ell +\lambda \left( \left( \omega \left( 1-\ell \right)
+rk_{0}\right) -c\left( 1+\tau \right) \right)
\end{equation*}%
so the focs are%
\begin{eqnarray*}
\frac{1}{c} &=&\lambda \left( 1+\tau \right) \\
\frac{1}{\ell } &=&\lambda \omega .
\end{eqnarray*}%
Note that this \textit{distortionary tax} shows up in the foc. This gives 
\begin{equation*}
\frac{\frac{1}{c}}{\frac{1}{\ell }}=\frac{1+\tau }{\omega }.
\end{equation*}
\end{frame}

\begin{frame}
\frametitle{The firm solves}
The firm solves:%
\begin{equation*}
\underset{k,n}{\max }k^{\alpha }n^{1-\alpha }-\omega n-rk
\end{equation*}%
which gives 
\begin{eqnarray*}
r &=&\alpha \left( \frac{n}{k}\right) ^{1-\alpha } \\
\omega &=&\left( 1-\alpha \right) \left( \frac{k}{n}\right) ^{\alpha }.
\end{eqnarray*}%
Compare (\ref{ce}) and (\ref{spp}) (repeated below).%
\begin{equation*}
\frac{\frac{1}{c}}{\frac{1}{\ell }}=\frac{1}{\left( 1-\alpha \right)
k^{\alpha }\left( 1-\ell \right) ^{-\alpha }}
\end{equation*}%
\begin{equation}
\frac{\frac{1}{c}}{\frac{1}{\ell }}=\frac{1+\tau }{\left( 1-\alpha \right)
k^{\alpha }\left( 1-\ell \right) ^{-\alpha }}.  \label{ce}
\end{equation}%
Thus the competitive equilibrium gives a different solution and is Pareto
inferior.
\end{frame}

\begin{frame}
\frametitle{Now let{'}s finish solving for the supply of labor in the CE. We have}
Now let's finish solving for the supply of labor in the CE. We have 
\begin{equation*}
\frac{\ell }{c}=\frac{1+\tau }{\left( 1-\alpha \right) k^{\alpha }\left(
1-\ell \right) ^{-\alpha }}.
\end{equation*}%
Putting in the budget constraint we get%
\begin{equation*}
\frac{\ell }{\left( 1+\tau \right) \left( \omega \left( 1-\ell \right)
+rk_{0}\right) }=\frac{1+\tau }{\left( 1-\alpha \right) k^{\alpha }\left(
1-\ell \right) ^{-\alpha }}
\end{equation*}%
or%
\begin{equation*}
\frac{\ell }{\left( \omega \left( 1-\ell \right) +rk_{0}\right) }=\frac{1}{%
\left( 1-\alpha \right) k^{\alpha }\left( 1-\ell \right) ^{-\alpha }}
\end{equation*}%
and using the equilibrium prices we get%
\begin{equation*}
\frac{\ell }{k^{\alpha }\left( 1-\ell \right) ^{1-\alpha }}=\frac{1}{\left(
1-\alpha \right) k^{\alpha }\left( 1-\ell \right) ^{-\alpha }}
\end{equation*}%
or%
\begin{equation}
\ell =\frac{1}{2-\alpha }.  \label{9}
\end{equation}%
Comparing with (\ref{6}) this is a different solution and Pareto
inferior.
\end{frame}

\begin{frame}
\frametitle{An economy with distortionary government expenditures}
Consider a similar representative agent economy where the consumer has a
utility function 
\begin{equation}
u=\ln c+\ln \ell  \label{u3}
\end{equation}%
where $c$ and $\ell $ are consumption and leisure. The consumer has an
endowment of 1 unit of time and $k_{0}$ units of capital. The representative
firm has a production technology given by $y=k^{\alpha }n^{1-\alpha }$ where 
$k$ and $n$ are capital and labor inputs. Suppose the government subsidizes
consumption. Specifically, for each unit of consumption goods purchased by
the representative consumer, she receives $s>0$ units of consumption from
the government. Thus letting $c_{p}$ be private spending on consumption, we
have $c=c_{p}\left( 1+s\right) .$ This subsidy is financed by a lump-sum tax 
$\tau $ on the representative consumer.

Find the solution to the social planner's problem.\newline
The social planner maximizes welfare subject to a \textit{resource constraint%
}. The problem is 
\begin{equation*}
\underset{c,\ell }{\max }\ln c+\ln \ell
\end{equation*}%
subject to:%
\begin{equation}
c=k^{\alpha }\left( 1-\ell \right) ^{1-\alpha }.  \notag
\end{equation}%
\end{frame}

\begin{frame}
\frametitle{Write this as}
Write this as%
\begin{equation*}
\underset{\ell }{\max }\ln k^{\alpha }\left( 1-\ell \right) ^{1-\alpha }+\ln
\ell .
\end{equation*}%
The foc\ is 
\begin{equation*}
\frac{1-\alpha }{1-\ell }=\frac{1}{\ell }.
\end{equation*}%
Solution is 
\begin{equation*}
\ell =\frac{1}{2-\alpha }.
\end{equation*}%
From this we can find $c$, etc.
\end{frame}

\begin{frame}
\frametitle{Define a competitive equilibrium for this economy}
Define a competitive equilibrium for this economy.\newline
A competitive equilibrium is a set of quantities $\left\{ y,c,\ell ,n,k,\tau
\right\} $ and prices $\left\{ \omega ,r\right\} $ such that

\begin{enumerate}
\item Each consumer chooses $c$ and $\ell $ to maximize (\ref{u3}) optimally
given $\omega ,r$,$s,$ and $\tau $.

\item The representative firm chooses $k$ and $n$ to maximize profits given $%
\omega $ and $r.$

\item The government choose $\tau $ to balance its budget.

\item Markets clear. That is:%
\begin{eqnarray*}
k &=&k_{0} \\
n+\ell &=&1 \\
c &=&y.
\end{eqnarray*}
\end{enumerate}
\end{frame}

\begin{frame}
\frametitle{Solve for the competitive equilibrium and show that it is not Pareto}
Solve for the competitive equilibrium and show that it is not Pareto
Optimal. The agent solves%
\begin{equation*}
\underset{c_{p},\ell }{\max }\ln c+\ln \ell
\end{equation*}%
subject to%
\begin{eqnarray*}
c &=&c_{p}\left( 1+s\right) \\
c_{p}+\tau &\leq &\omega \left( 1-\ell \right) +rk_{0} \\
\ell &\in &\left[ 0,1\right] .
\end{eqnarray*}%
\end{frame}

\begin{frame}
\frametitle{As the second and third hold with equality in equilibrium, write this as}
As the second and third hold with equality in equilibrium, write this as%
\begin{equation*}
\underset{c_{p},\ell }{\max }\ln c_{p}\left( 1+s\right) +\ln \ell +\lambda
\left( \omega \left( 1-\ell \right) +rk_{0}-c_{p}-\tau \leq \right)
\end{equation*}%
and focs give 
\begin{equation*}
\frac{\frac{1}{c_{p}}}{\frac{1}{\ell }}=\frac{1}{\omega }
\end{equation*}%
or%
\begin{equation*}
\frac{\frac{c}{1+s}}{\frac{1}{\ell }}=\frac{1}{\omega }
\end{equation*}%
The subsidy is distoring the tradeoff beteen consumption and leisure.
\end{frame}

\begin{frame}
\frametitle{We would be more likely the write this as}
We would be more likely the write this as 
\begin{equation*}
\underset{\ell }{\max }\ln \left[ \left( 1+s\right) \left( \omega \left(
1-\ell \right) +rk_{0}-\tau \right) \right] +\ln \ell .
\end{equation*}%
The foc is 
\begin{equation}
\frac{\omega }{\left( \omega \left( 1-\ell \right) +rk_{0}-\tau \right) }=%
\frac{1}{\ell }.  \label{nfoc}
\end{equation}%
\end{frame}

\begin{frame}
\frametitle{The firm solves}
The firm solves%
\begin{equation*}
\underset{k,n}{\max }k^{\alpha }n^{1-\alpha }-\omega n-rk
\end{equation*}%
which gives 
\begin{eqnarray}
r &=&\alpha \left( \frac{n}{k}\right) ^{1-\alpha }  \label{wages} \\
\omega &=&\left( 1-\alpha \right) \left( \frac{k}{n}\right) ^{\alpha }. 
\notag
\end{eqnarray}%
The government budget constraint is%
\begin{equation*}
c_{p}s=\tau
\end{equation*}%
or 
\begin{equation}
\frac{cs}{1+s}=\tau .  \label{bud2}
\end{equation}%
\end{frame}

\begin{frame}
\frametitle{Slide 41}
Putting (\ref{wages}), (\ref{bud2}), $k=k_{0}$ and $n=1-\ell $ and $c=y$ and 
$y=k^{\alpha }n^{1-\alpha }$ into (\ref{nfoc}) gives 
\begin{equation*}
\frac{\left( 1-\alpha \right) \left( \frac{k}{1-\ell }\right) ^{\alpha }}{%
\left( 1-\alpha \right) \left( \frac{k}{1-\ell }\right) ^{\alpha }\left(
1-\ell \right) +\alpha \left( \frac{1-\ell }{k}\right) ^{1-\alpha }k-\frac{%
sk^{\alpha }\left( 1-\ell \right) ^{1-\alpha }}{1+s}}=\frac{1}{\ell }
\end{equation*}%
which simplifies to 
\begin{equation*}
\frac{\left( 1-\alpha \right) }{\left( 1-\alpha \right) \left( 1-\ell
\right) +\alpha \left( 1-\ell \right) -\frac{s\left( 1-\ell \right) }{1+s}}=%
\frac{1}{\ell },
\end{equation*}%
\begin{equation*}
\frac{\left( 1+s\right) \left( 1-\alpha \right) }{1-\ell }=\frac{1}{\ell }.
\end{equation*}%
The solution is 
\begin{equation*}
\ell =\frac{1}{2-\alpha +s\left( 1-\alpha \right) }
\end{equation*}%
so clearly any positive (or negative) subsidy will be inefficient.
\end{frame}

\begin{frame}
\frametitle{An economy with a missing market: valued government expenditures}
Consider a similar representative agent economy where the consumer has a
utility function 
\begin{equation*}
u=\ln c+\ln \ell +\ln g
\end{equation*}%
where $c$ and $\ell $ are consumption and leisure and $g$ is the supply of
goods which are valued by agents and provided by government. Note that since
we assume only government can provide this good, we can consider this a case
of \textit{missing markets} for the good $g$. The consumer has an endowment
of 1 unit of time. The representative firm has a production technology given
by $y=n$ where $n$ is the labor input. Government spending is funded by a
lump-sum tax $\tau $ on the representative consumer. A unit of output
converts one-for-one to a unit of $g$. The social planner maximizes welfare
subject to a \textit{resource constraint}. It solves 
\begin{equation*}
\underset{\ell ,c,g}{\max }\ln c+\ln \ell +\ln g
\end{equation*}%
subject to%
\begin{equation}
c+g=1-\ell .  \notag
\end{equation}%
\end{frame}

\begin{frame}
\frametitle{Write this as}
Write this as%
\begin{equation*}
\underset{\ell ,g}{\max }\ln \left( 1-\ell -g\right) +\ln \ell +\ln g.
\end{equation*}%
The focs\ are 
\begin{equation*}
\frac{1}{1-\ell -g}=\frac{1}{\ell }
\end{equation*}%
\begin{equation*}
g=\frac{1-\ell }{2}.
\end{equation*}%
\end{frame}

\begin{frame}
\frametitle{This into the first gives}
This into the first gives%
\begin{equation*}
\frac{1}{1-\ell -\frac{1-\ell }{2}}=\frac{1}{\ell }
\end{equation*}%
\begin{equation*}
\ell =\frac{1}{3}.
\end{equation*}%
so 
\begin{equation*}
g=\frac{1-\frac{1}{3}}{2}=\frac{1}{3}
\end{equation*}%
\begin{equation*}
c=\frac{1-\frac{1}{3}}{2}=\frac{1}{3}.
\end{equation*}
\end{frame}

\begin{frame}
\frametitle{Define a competitive equilibrium for this economy}
Define a competitive equilibrium for this economy.\newline
A competitive equilibrium is a set of quantities $\left\{ y,c,\ell ,n,\tau
\right\} $ and price $\left\{ \omega \right\} $ such that

\begin{enumerate}
\item Each consumer chooses $c$ and $\ell $ to maximize (\ref{u3}) optimally
given $\omega ,$ $g$ and $\tau $.

\item The representative firm chooses $n$ to maximize profits given $\omega
. $

\item The government balances its budget.

\item Markets clear. That is:%
\begin{eqnarray*}
n+\ell &=&1 \\
c+g &=&y.
\end{eqnarray*}
\end{enumerate}
\end{frame}

\begin{frame}
\frametitle{Now solve for the competitive equilibrium. The agent solves}
Now solve for the competitive equilibrium. The agent solves:%
\begin{equation*}
\underset{c,\ell }{\max }\ln c+\ln \ell +\ln g
\end{equation*}%
subject to%
\begin{eqnarray*}
c &\leq &\left( \omega \left( 1-\ell \right) -\tau \right) \\
\ell &\in &\left[ 0,1\right]
\end{eqnarray*}%
where $g$ is taken as given by the agent. As the first holds with equality
in equilibrium write this as 
\begin{equation*}
\underset{\ell }{\max }\ln \left[ \omega \left( 1-\ell \right) -\tau \right]
+\ln \ell +\ln g.
\end{equation*}%
The foc is 
\begin{equation}
\frac{\omega }{\omega \left( 1-\ell \right) -\tau }=\frac{1}{\ell }.
\label{foc4}
\end{equation}%
\end{frame}

\begin{frame}
\frametitle{The firm{'} problem gives}
The firm' problem gives 
\begin{equation}
\omega =1.  \notag
\end{equation}%
The government budget constraint is%
\begin{equation}
g=\tau .  \label{bc}
\end{equation}%
Putting $\omega =1$ into (\ref{foc4}) gives 
\begin{eqnarray}
\frac{1}{\left( 1-\ell \right) -\tau } &=&\frac{1}{\ell }  \label{foc5} \\
\ell &=&\frac{1-\tau }{2}.  \notag
\end{eqnarray}%
$\allowbreak $How should the government choose $g?$
\end{frame}

\begin{frame}
\frametitle{Slide 48}
We want 
\begin{equation*}
\ell =\frac{1}{3}
\end{equation*}%
From above%
\begin{eqnarray*}
\frac{1}{3} &=&\frac{1-\tau }{2} \\
\tau &=&\frac{1}{3}.
\end{eqnarray*}%
From government budget constraint%
\begin{equation*}
g=\frac{1}{3}.
\end{equation*}%
This is the same as we found using the social planner's problem. Here we
have shown this is consistent with the proper supply of labor hours as well.
Thus we have shown that through proper policy the social optimum can be
achieved and that the social planner's problem can be used to find the
optimal policy.
\end{frame}

\begin{frame}
\frametitle{An economy with an externality}
Here government has a role even though expenditures are not directly valued.
Government can overcome an externality. Consider an economy with many
agents, each of whom maximizes 
\begin{equation*}
\frac{c^{\gamma }}{\gamma }.
\end{equation*}%
Here $c$ is consumption and $\gamma <1.$ Each agent has one unit of time
available in each period and must accumulate human capital before he works.
An individual accumulates human capital according to $h=\left( 1-u\right) $
where $u$ is the share of time spent working$.$ A representative firm
produces according to $y=\alpha H$ where $H$ is the amount of human capital
employed and equal to $hu$ in equilibrium. The social planner's problem is $%
\ $%
\begin{equation*}
\underset{c,u,h}{\max }\frac{c^{\gamma }}{\gamma }
\end{equation*}%
subject to 
\begin{eqnarray*}
c &=&\alpha H \\
H &=&\left( 1-u\right) u.
\end{eqnarray*}%
is the amount of human capital provided per unit of time, and $u$ in the
share of time spent in production.
\end{frame}

\begin{frame}
\frametitle{Rewrite as}
Rewrite as%
\begin{equation*}
\underset{u}{\max }\frac{\left( \alpha (1-u)u\right) ^{\gamma }}{\gamma }.
\end{equation*}%
It is easy to show $u^{\ast }=1/2$.\ \newline
Now suppose instead that human capital is produced according to $h=\overline{%
h}^{1-\theta }\left( 1-u\right) ^{\theta }$ where $0<\theta <1$ and $%
\overline{h}$ is the average human capital across all agents in the economy.
There is an externality in human capital accumulation. The representative
firm still produces output accordind to $y=\alpha H.$ This does not change
the social planner's problem. Why? Now look at CE. The agent's problem is%
\begin{equation*}
\underset{c,u,h}{\max }\frac{c^{\gamma }}{\gamma }
\end{equation*}%
subject to 
\begin{equation*}
\begin{array}{ll}
c=whu; & h=\overline{h}^{1-\theta }\left( 1-u\right) ^{\theta }%
\end{array}%
\end{equation*}%
where $w$ is the wage\ per unit of human and $w=\alpha $ in
equilibrium.
\end{frame}

\begin{frame}
\frametitle{Rewrite this as}
Rewrite this as%
\begin{equation*}
\underset{u}{\max }\frac{\left( whu\right) ^{\gamma }}{\gamma }.
\end{equation*}%
or%
\begin{equation*}
\underset{u}{\max }\frac{\left( w\overline{h}^{1-\theta }(1-u)^{\theta
}u\right) ^{\gamma }}{\gamma }.
\end{equation*}%
This is maximized at 
\begin{equation*}
u=\frac{1}{1+\theta }.
\end{equation*}%
Notice that the social planner's solution and the result of a competitive
equilibrium do not coincide. Why?
\end{frame}

\begin{frame}
\frametitle{Now suppose that the government pays people to go to school and...}
Now suppose that the government pays people to go to school and taxes them
lump sum to pay for the subsidies. Specifically the government pays $%
\varsigma w$ per unit of time in school where $\varsigma $ is the rate at
which schooling is subsidized$.$ The agent's problem is 
\begin{equation*}
\underset{u}{\max }\frac{\left( whu+\varsigma w(1-u)-\tau \right) ^{\gamma }%
}{\gamma }.
\end{equation*}%
Write this as%
\begin{equation*}
\underset{u}{\max }\frac{\left( w\overline{h}^{1-\theta }(1-u)^{\theta
}u+\varsigma w(1-u)-\tau \right) ^{\gamma }}{\gamma }
\end{equation*}%
which is maximized where $\overline{h}^{1-\theta }(1-u)^{\theta }u+\varsigma
(1-u)$ is maximized.
\end{frame}

\begin{frame}
\frametitle{This gives}
This gives%
\begin{equation*}
-\theta \overline{h}^{1-\theta }(1-u)^{\theta -1}u+\overline{h}^{1-\theta
}(1-u)^{\theta }-\varsigma =0
\end{equation*}%
In equilibrium $\overline{h}=h=(1-u)$ so 
\begin{equation*}
-\theta \left( 1-u\right) ^{1-\theta }(1-u)^{\theta -1}u+(1-u)^{1-\theta
}(1-u)^{\theta }-\varsigma =0
\end{equation*}%
\begin{equation*}
-\theta u+(1-u)-\varsigma =0
\end{equation*}%
\begin{equation*}
u=\frac{1-\varsigma }{\theta +1}.
\end{equation*}%
$\allowbreak \allowbreak $To get the social planner's solution, government
should set $\varsigma $ such that 
\begin{equation*}
\frac{1}{2}=\frac{1-\varsigma }{\theta +1}.
\end{equation*}%
Thus $\varsigma =\frac{1-\theta }{2}.$
\end{frame}

\begin{frame}
\frametitle{An economy with imperfect competition}
Consider an economy with a single final good which is produced using a
variety of intermediate goods. Specifically, $n$ intermediate goods ($x_{i},$
$i\in \left\{ 1,2..n\right\} $) combine with labor ($L$) to produce a final
output good $y$ according to 
\begin{equation*}
y=AL^{1-\alpha }\overset{n}{\underset{i=1}{\sum }}x_{i}^{\alpha }
\end{equation*}%
where $A>0$ and $0<\alpha <1.$ The final good is the numeraire good and is
produced competitively. There are many producers, but they are identical.
For simplicity we consider a single representative firm. For a producer of
the final good, the firm's profits $\left( \pi _{f}\right) $ are given by 
\begin{equation*}
\pi _{f}=y-\overset{n}{\underset{i=1}{\sum }p_{i}}x_{i}-wL
\end{equation*}%
where $p_{i}$ is the relative price of intermediate good $i$ and $w$ is the
wage. Profits can be rewritten as%
\begin{equation*}
\pi _{f}=AL^{1-\alpha }\overset{n}{\underset{i=1}{\sum }}x_{i}^{\alpha }-%
\overset{n}{\underset{i=1}{\sum }p_{i}}x_{i}-wL.
\end{equation*}%
\end{frame}

\begin{frame}
\frametitle{In choosing the foc for the firm{'}s optimization problem is}
In choosing $x_{i},$ the foc for the firm's optimization problem is%
\begin{equation*}
\alpha AL^{1-\alpha }x_{i}^{\alpha -1}=p_{i}
\end{equation*}%
from which we can get the \textit{derived demand} for intermediate good $%
x_{i}.$ 
\begin{equation*}
x_{i}=\left( \frac{\alpha AL^{1-\alpha }}{p_{i}}\right) ^{\frac{1}{1-\alpha }%
}.
\end{equation*}%
Each intermediate good is unique and produced by only one firm. Thus there
is \textit{imperfect competition} and the producer of an intermediate good
is not a price taker. The only input to the production of an intermediate
good is one unit of the final good. Thus the profit for the intermediate
firm is 
\begin{equation*}
\pi _{i}=p_{i}x_{i}-x_{i}
\end{equation*}%
or 
\begin{equation*}
\pi _{i}=\left( p_{i}-1\right) \left( \frac{\alpha AL^{1-\alpha }}{p_{i}}%
\right) ^{\frac{1}{1-\alpha }}.
\end{equation*}
\end{frame}

\begin{frame}
\frametitle{The firm chooses price to maximize profit. The foc is}
The firm chooses price to maximize profit. The foc is%
\begin{equation*}
\left( p_{i}-1\right) \frac{1}{1-\alpha }\left( \frac{\alpha AL^{1-\alpha }}{%
p_{i}}\right) ^{\frac{1}{1-\alpha }-1}\frac{\alpha AL^{1-\alpha }}{p_{i}^{2}}%
=\left( \frac{\alpha AL^{1-\alpha }}{p_{i}}\right) ^{\frac{1}{1-\alpha }}
\end{equation*}%
or%
\begin{equation*}
\left( p_{i}-1\right) \frac{1}{1-\alpha }\frac{1}{p_{i}^{2}}=\left( \frac{1}{%
p_{i}}\right)
\end{equation*}%
so 
\begin{equation*}
p_{i}=\frac{1}{\alpha }.
\end{equation*}%
Interpret this.\textit{\ }
\end{frame}

\begin{frame}
\frametitle{Slide 57}
\ 

Thus%
\begin{equation*}
x_{i}=\left( \alpha ^{2}AL^{1-\alpha }\right) ^{\frac{1}{1-\alpha }}.
\end{equation*}%
Notice that it is the same as for all goods so that%
\begin{equation*}
y=AL^{1-\alpha }\overset{n}{\underset{i=1}{\sum }}x_{i}^{\alpha }
\end{equation*}%
becomes%
\begin{equation*}
y=nAL^{1-\alpha }\left( \alpha ^{2}AL^{1-\alpha }\right) ^{\frac{\alpha }{%
1-\alpha }}
\end{equation*}%
and goods available for consumption are 
\begin{equation*}
\overset{\text{output}}{\overbrace{nAL^{1-\alpha }\left( \alpha
^{2}AL^{1-\alpha }\right) ^{\frac{\alpha }{1-\alpha }}}}-\overset{\text{used
to produce inter. goods}}{\overbrace{n\left( \alpha ^{2}AL^{1-\alpha
}\right) ^{\frac{1}{1-\alpha }}}}.
\end{equation*}%
$\allowbreak $
\end{frame}

\begin{frame}
\frametitle{We can show that the final goods producer has zero profits}
\begin{itemize}
\item We can \ show that the final goods producer has zero profits.

\item Does the producer of the intermediate good have profits? Rewrite%
\begin{equation*}
\pi _{i}=p_{i}x_{i}-x_{i}
\end{equation*}%
as%
\begin{equation*}
\pi _{i}=\left( \alpha ^{2}AL^{1-\alpha }\right) ^{\frac{1}{1-\alpha }%
}\left( \frac{1}{\alpha }-1\right) .
\end{equation*}%
Obviously there are profits unless $\alpha =1.$

\item To find income add labor income to dividend income (profits).

\item We can show the goods market clears.

\item Preferences? Since there is no labor leisure choice and only one good
and only one time period, the agent simply buys goods with the proceeds of
her entire endowment. Any set of preferences will give us the same result so
long as utility increases in consumption.

\item Is the competitive equilibrium Pareto optimal? Here maximizing
consumption is the same as maximizing utility. So the question whether the
competitive equilibrium brings about the greatest possible consumption.
\end{itemize}
\end{frame}

\begin{frame}
\frametitle{Consider the social planner{'}s problem}
Consider the social planner's problem. 
\begin{equation*}
\underset{x_{i}}{\max }\left( AL^{1-\alpha }\overset{n}{\underset{i=1}{\sum }%
}x_{i}^{\alpha }-\overset{n}{\underset{i=1}{\sum }}x_{i}\right) .
\end{equation*}%
The foc is 
\begin{equation*}
\alpha AL^{1-\alpha }x_{i}^{\alpha -1}=1
\end{equation*}%
\begin{equation*}
x_{i}^{\ast }=\left( \alpha AL^{1-\alpha }\right) ^{\frac{1}{1-\alpha }}.
\end{equation*}%
Compare this to 
\begin{equation*}
x_{i}=\left( \alpha ^{2}AL^{1-\alpha }\right) ^{\frac{1}{1-\alpha }}.
\end{equation*}%
Clearly, $x_{i}<x_{i}^{\ast }$. Why?
\end{frame}

\begin{frame}
\frametitle{Can policy help? Suppose we decide to lump sum tax all individuals...}
Can policy help? Suppose we decide to lump sum tax all individuals to fund a
production subsidy. Each time an intermediate firm produces, it is given $%
\varsigma $ units of the final good as a subsidy 
\begin{equation*}
\pi _{i}=p_{i}x_{i}-x_{i}+\varsigma x_{i}
\end{equation*}%
\begin{equation*}
\pi _{i}=p_{i}x_{i}-x_{i}\left( 1-\varsigma \right) .
\end{equation*}%
The we could show 
\begin{equation*}
p_{i}=\frac{\left( 1-\varsigma \right) }{\alpha }.
\end{equation*}%
Recall that the demand for the intermediate good is 
\begin{equation*}
x_{i}=\left( \frac{\alpha AL^{1-\alpha }}{p_{i}}\right) ^{\frac{1}{1-\alpha }%
}
\end{equation*}%
or%
\begin{equation*}
x_{i}=\left( \frac{\alpha AL^{1-\alpha }}{\frac{\left( 1-\varsigma \right) }{%
\alpha }}\right) ^{\frac{1}{1-\alpha }}.
\end{equation*}%
The optimum can be achieved if $\varsigma =1-\alpha .$ Discuss how such a
model might give insights into the growth process (R\&D).
\end{frame}

\begin{frame}
\frametitle{An economy with simple dynamics}
We now introduce \textit{dynamics} in a simple 2 period exchange economy. We
assume no uncertainty, competitive markets and non-storable goods.

\textit{Preferences: }an agent has log utility over consumption in the first
and second periods of life ($c_{1}$ and $c_{2}$), that is 
\begin{equation*}
u\left( c_{1},c_{2}\right) =\ln c_{1}+\beta \ln c_{2}
\end{equation*}%
where $0<\beta <1.$

\textit{Endowments:} the representative agent is endowed with quantities $%
y_{1}$ and $y_{2}$ respectively in the first and second period. The agent
maximizes utility subject to%
\begin{equation*}
p_{1}c_{1}+p_{2}c_{2}\leq p_{1}y_{1}+p_{2}y_{2}
\end{equation*}%
where prices are market clearing prices.

A competitive equilibrium in this economy is a set of prices $\left\{
p_{1},p_{2}\right\} $ and allocations $\left\{ c_{1},c_{2}\right\} $ such
that:

\begin{enumerate}
\item The consumer maximizes utility subject to the budget constraint for
given prices.

\item Markets clear: $c_{t}=y_{t}$ for $t=1,2$.
\end{enumerate}
\end{frame}

\begin{frame}
\frametitle{Example. Let}
\textit{Example}. Let%
\begin{equation*}
u\left( c_{1},c_{2}\right) =\ln c_{1}+\beta \ln c_{2}+\lambda \left(
p_{1}y_{1}+p_{2}y_{2}-p_{1}c_{1}-p_{2}c_{2}\right) .
\end{equation*}%
Focs are:%
\begin{equation*}
\frac{1}{c_{1}}=\lambda p_{1}
\end{equation*}%
\begin{equation*}
\frac{\beta }{c_{2}}=\lambda p_{2}.
\end{equation*}%
These imply $c_{2}=\frac{p_{1}}{p_{2}}\beta c_{1}.$ This into $%
p_{1}c_{1}+p_{2}c_{2}=p_{1}y_{1}+p_{2}y_{2}$ gives the following demand
functions%
\begin{equation*}
c_{1}=\frac{1}{1+\beta }\left( y_{1}+\frac{p_{2}}{p_{1}}y_{2}\right)
\end{equation*}%
\begin{equation*}
c_{2}=\frac{\beta }{1+\beta }\frac{p_{1}}{p_{2}}\left( y_{1}+\frac{p_{2}}{%
p_{1}}y_{2}\right) .
\end{equation*}%
\end{frame}

\begin{frame}
\frametitle{Slide 63}
We now impose the second half of our definition of an equilibrium. $%
c_{1}=y_{1}$ and $c_{2}=y_{2}$. We need only impose one of these (Walras'
Law). Using $c_{1}=y_{1}$ gives 
\begin{equation*}
y_{1}=\frac{1}{1+\beta }\left( y_{1}+\frac{p_{2}}{p_{1}}y_{2}\right)
\end{equation*}%
and solving for prices gives:%
\begin{equation*}
\frac{p_{2}}{p_{1}}=\beta \frac{y_{1}}{y_{2}}
\end{equation*}%
These prices reflect the rate that an agent \textit{would} trade current for
future consumption. We will soon allow such trade. Note: We can only solve
for relative prices $\frac{p_{1}}{p_{2}}$. We can set $p_{1}=1$ without
losing any generality.

Thus the prices $\left\{ p_{1}=1,\text{ }p_{2}=\beta \frac{y_{1}}{y_{2}}%
\right\} $ and allocations $\left\{ c_{1}=y_{1},\text{ }c_{2}=y_{2}\right\} $
constitute a competitive equilibrium for this economy.

Note this is 
\begin{equation*}
\frac{p_{1}}{p_{2}}=\frac{1}{p_{2}}=\frac{\frac{1}{c_{1}}}{\frac{\beta }{%
c_{2}}}=\frac{u_{1}}{\beta u_{2}}=\left\vert MRS_{1,2}\right\vert
\end{equation*}%
With output unable to adjust, the price of a period 2 good adjusts so that
the ratio of prices is equal to the absolute value of the marginal rate of
substitution.
\end{frame}

\begin{frame}
\frametitle{An economy with savings}
Suppose there is a bond market. The price of a bond is one unit of the
consumption good. Each bond purchased in period 1 yields $\left(
1+r_{b}\right) $ units of the consumption good in period two. The budget
constraint in periods 1 and two are%
\begin{eqnarray*}
p_{1}c_{1}+p_{1}b &=&p_{1}y_{1} \\
c_{2}p_{2} &=&p_{2}b\left( 1+r_{b}\right) +p_{2}y_{2}
\end{eqnarray*}%
or%
\begin{eqnarray*}
c_{1}+b &=&y_{1} \\
c_{2} &=&b\left( 1+r_{b}\right) +y_{2}.
\end{eqnarray*}%
\end{frame}

\begin{frame}
\frametitle{We can write the Lagrangian as}
We can write the Lagrangian as 
\begin{equation*}
L=\ln c_{1}+\beta \ln c_{2}+\lambda _{1}\left( y_{1}-b-c_{1}\right) +\lambda
_{1}\left( y_{2}+b\left( 1+r_{b}\right) -c_{2}\right) .
\end{equation*}

Focs are 
\begin{eqnarray*}
\frac{1}{c_{1}} &=&\lambda _{1} \\
\frac{\beta }{c_{2}} &=&\lambda _{2} \\
\lambda _{1} &=&\lambda _{2}\left( 1+r_{b}\right) \\
c_{1} &=&y_{1}-b \\
c_{2} &=&y_{2}+b\left( 1+r_{b}\right) .
\end{eqnarray*}%
\end{frame}

\begin{frame}
\frametitle{In equilibrium, the interest rate must adjust so that the goods market}
\begin{itemize}
\item In equilibrium, the interest rate must adjust so that the goods market
clears in each period.

\item This is the same role played by prices in the earlier statement of the
model. A competitive equilibrium in this economy is a price $\left\{
r_{b}\right\} $ and allocations $\left\{ c_{1},c_{2}\right\} $ such that:
\end{itemize}

\begin{enumerate}
\item The consumer maximizes utility subject to the budget constraint for
given prices.

\item Markets clear: $c_{t}=y_{t}$ for $t=1,2$ and $b=0$. (Why $b=0$)?
\end{enumerate}

From the first three focs we have%
\begin{equation*}
c_{2}=\beta c_{1}\left( 1+r_{b}\right)
\end{equation*}%
Then from the two market clearing conditions we have%
\begin{equation*}
y_{2}=\beta y_{1}\left( 1+r_{b}\right)
\end{equation*}%
and the equilibrium interest rate is%
\begin{equation*}
1+r_{b}=\frac{y_{2}}{\beta y_{1}}
\end{equation*}%
\begin{equation*}
\frac{1}{1+r_{b}}=\frac{\beta y_{1}}{y_{2}}
\end{equation*}%
\end{frame}

\begin{frame}
\frametitle{Slide 67}
Recall%
\begin{equation*}
p_{2}=\frac{\beta y_{1}}{y_{2}}
\end{equation*}%
Compare this to our earlier expression to find that $\frac{1}{p_{2}}=\left(
1+r_{b}\right) \ $or $p_{2}=\frac{1}{1+r_{b}}.$ The price of a future good
is the discounted value of a current good, since one current good can be
transformed into $1+r\ $future goods.

Note this is 
\begin{equation*}
\frac{p_{1}}{p_{2}}=\frac{1}{p_{2}}=\left( 1+r_{b}\right) =\frac{\frac{1}{%
c_{1}}}{\frac{\beta }{c_{2}}}=\frac{u_{1}}{\beta u_{2}}=\left\vert
MRS_{1,2}\right\vert
\end{equation*}%
With output unable to adjust, the interest rate adjusts so that $\left(
1+r_{b}\right) $ is equal to the absolute value of the marginal rate of
substitution. Prices adjust accordingly, but knowing $(1+r_{b})$ is
sufficient.
\end{frame}

\begin{frame}
\frametitle{Two other common approaches}
Two other common approaches.

The first other common approach is to combine period budget constraints into
a lifetime budget constraint. Recall 
\begin{eqnarray*}
c_{1}+b &=&y_{1} \\
c_{2} &=&b\left( 1+r_{b}\right) +y_{2}.
\end{eqnarray*}%
The second of these gives 
\begin{equation*}
\frac{c_{2}-y_{2}}{1+r_{b}}=b
\end{equation*}%
and this into the first gives%
\begin{eqnarray*}
c_{1}+\frac{c_{2}-y_{2}}{1+r_{b}} &=&y_{1} \\
c_{1}+\frac{c_{2}}{1+r_{b}} &=&y_{1}+\frac{y_{2}}{1+r_{b}}.
\end{eqnarray*}%
\end{frame}

\begin{frame}
\frametitle{We can write the problem as}
We can write the problem as 
\begin{equation*}
L=\ln c_{1}+\beta \ln c_{2}+\lambda \left( y_{1}+\frac{c_{2}}{1+r_{b}}-c_{1}-%
\frac{c_{2}}{1+r_{b}}\right)
\end{equation*}%
so that focs are%
\begin{eqnarray*}
\frac{1}{c_{1}} &=&\lambda \\
\frac{\beta }{c_{2}} &=&\frac{\lambda }{1+r_{b}} \\
y_{1}+\frac{c_{2}}{1+r_{b}} &=&c_{1}+\frac{c_{2}}{1+r_{b}}.
\end{eqnarray*}%
\end{frame}

\begin{frame}
\frametitle{Together these give}
Together these give 
\begin{equation*}
c_{2}=\beta c_{1}\left( 1+r_{b}\right) .
\end{equation*}%
This into the budget constraint gives%
\begin{equation*}
c_{1}+\beta c_{1}=y_{1}+\frac{y_{2}}{1+r_{b}}.
\end{equation*}%
In equilibrium (using goods market clearing) 
\begin{equation*}
y_{1}+\beta y_{1}=y_{1}+\frac{y_{2}}{1+r_{b}}
\end{equation*}%
\begin{equation*}
1+r_{b}=\frac{y_{2}}{\beta y_{1}}.
\end{equation*}%
\end{frame}

\begin{frame}
\frametitle{The second other common approach is to plug in period budget...}
The second other common approach is to plug in period budget constraints. We
can write the problem as%
\begin{equation*}
L=\ln \left( y_{1}-b\right) +\beta \ln \left( y_{2}+b\left( 1+r_{b}\right)
\right) .
\end{equation*}%
so foc gives 
\begin{equation*}
\frac{1}{y_{1}-b}=\frac{\beta \left( 1+r_{b}\right) }{y_{2}+\left(
1+r_{b}\right) b}.
\end{equation*}%
Now impose $b=0$ in equilibrium giving 
\begin{equation*}
\frac{1}{y_{1}}=\frac{\beta \left( 1+r_{b}\right) }{y_{2}}
\end{equation*}%
or 
\begin{equation*}
1+r_{b}=\frac{y_{2}}{\beta y_{1}}.
\end{equation*}%
\end{frame}

\begin{frame}
\frametitle{An economy with storage}
Above there is no opportunity for producing goods and $r$ depends on
preferences and endowments. In general $r$ will also depend on the
technological opportunities for transforming current consumption to future
consumption. Consider the case above except that there is a storage
technology: one good can be transformed to $1+r_{k}\;$units tomorrow. A bond
market still exists. Let $k$ be the amount of storage chosen by the
representative agent. We leave out prices as indicated above. The first and
second period budget constraints are 
\begin{eqnarray*}
c_{1}+b+k &=&y_{1} \\
c_{2} &=&k\left( 1+r_{k}\right) +b\left( 1+r_{b}\right) +y_{2}.
\end{eqnarray*}%
The individual's problem is%
\begin{equation*}
L=\ln c_{1}+\beta \ln c_{2}+\lambda _{1}\left( y_{1}-c_{1}-k-b\right)
+\lambda _{2}\left( y_{2}+k\left( 1+r_{k}\right) +b\left( 1+r_{b}\right)
-c_{2}\right)
\end{equation*}%
and focs are 
\begin{eqnarray*}
\frac{1}{c_{1}} &=&\lambda _{1} \\
\frac{\beta }{c_{2}} &=&\lambda _{2} \\
\lambda _{1} &=&\lambda _{2}\left( 1+r_{b}\right) \\
\lambda _{1} &=&\lambda _{2}\left( 1+r_{k}\right) \\
y_{1} &=&c_{1}+k+b \\
y_{2}+k\left( 1+r_{k}\right) +b\left( 1+r_{b}\right) &=&c_{2}.
\end{eqnarray*}
\end{frame}

\begin{frame}
\frametitle{So these imply or (assuming for now all constraints}
So these imply $\left( 1+r_{k}\right) =\left( 1+r_{b}\right) $ or $\left(
1+r_{k}\right) =\left( 1+r_{b}\right) =r.$ (assuming for now all constraints
hold with equality)

\begin{itemize}
\item In our earlier example, $\left( 1+r_{b}\right) $ adjusted so that
consumption equaled endowments.

\item Here consumption is adjusting so that $\left( 1+r_{k}\right) =\left(
1+r_{b}\right) .\ $Let's call this common return $R$.

\item Bond market clearing still gives $b=0.$

\item Goods market clearing conditions are now 
\begin{eqnarray*}
c_{1} &=&y_{1}-k \\
c_{2} &=&y_{2}+kR.
\end{eqnarray*}%
\end{itemize}
\end{frame}

\begin{frame}
\frametitle{Slide 74}
From focs 
\begin{equation*}
c_{2}=\beta c_{1}R
\end{equation*}%
Note this is 
\begin{equation*}
\frac{1}{p_{2}}=\left( 1+r_{b}\right) =\left( 1+r_{k}\right) =R=\frac{\frac{1%
}{c_{1}}}{\frac{\beta }{c_{2}}}=\frac{u_{1}}{\beta u_{2}}=\left\vert
MRS_{1,2}\right\vert
\end{equation*}%
With prices and $R$ unable to adjust, consumption adjusts through savings so
that $R$ is equal to the absolute value of the marginal rate of
substitution.
\end{frame}

\begin{frame}
\frametitle{Before, we were done at this point since was equal to . Now}
Before, we were done at this point since $c_{i}$ was equal to $y_{i}$. Now
we need to go a bit further. From focs 
\begin{equation*}
c_{2}=\beta c_{1}R
\end{equation*}%
so from second goods market clearing condition%
\begin{equation*}
\beta c_{1}R=y_{2}+kR
\end{equation*}%
or 
\begin{equation*}
k=\beta c_{1}-\frac{y_{2}}{R}
\end{equation*}%
and this into the first gives 
\begin{equation*}
c_{1}+\beta c_{1}-\frac{y_{2}}{R}=y_{1}
\end{equation*}%
or 
\begin{equation*}
c_{1}=\frac{1}{1+\beta }\left( y_{1}+\frac{y_{2}}{R}\right) .
\end{equation*}%
\end{frame}

\begin{frame}
\frametitle{Note 1: In the problem above, savings for the individual was . We}
\end{itemize}

Note 1: In the problem above, savings for the individual was $s=k+b$. We
again define $R$ to be general returns on saving and write 
\begin{equation*}
L=\ln c_{1}+\beta \ln c_{2}+\lambda _{1}\left( y_{1}-c_{1}-s\right) +\lambda
_{2}\left( y_{2}+sR-c_{2}\right) .
\end{equation*}%
or combine period constraints to a lifetime constraint and write 
\begin{equation*}
L=\ln c_{1}+\beta \ln c_{2}+\lambda \left( y_{1}+\frac{y_{2}}{R}-c_{1}-\frac{%
c_{2}}{R}\right)
\end{equation*}%
or plug in period budget constraints and write%
\begin{equation*}
L=\ln \left( y_{1}-s\right) +\beta \ln \left( y_{2}+sR\right) .
\end{equation*}%
Then we solve for savings and use no arbitrage conditions to argue $%
R=1+r_{k}=1+r_{b}.$ Also in equilibrium with $b=0$ we have $s=k+b=k$.\bigskip
\end{frame}

\begin{frame}
\frametitle{Note 2: It is common also to solve only for and let the bond market}
Note 2: It is common also to solve only for $k$ and let the bond market
operate in the background. Then we just write the problem as%
\begin{equation*}
L=\ln c_{1}+\beta \ln c_{2}+\lambda _{1}\left( y_{1}-c_{1}-k\right) +\lambda
_{2}\left( y_{2}+k\left( 1+r_{k}\right) -c_{2}\right)
\end{equation*}%
or combine period constraints to a lifetime constraint and write 
\begin{equation*}
L=\ln c_{1}+\beta \ln c_{2}+\lambda \left( y_{1}+\frac{y_{2}}{1+r_{k}}-c_{1}-%
\frac{c_{2}}{1+r_{k}}\right)
\end{equation*}%
$\left( y_{1}-c_{1}-\frac{c_{2}-y_{2}}{r_{k}+1}\right) $or plug in period
budget constraints and write%
\begin{equation*}
L=\ln \left( y_{1}-k\right) +\beta \ln \left( y_{2}+k\left( 1+r_{k}\right)
\right) .
\end{equation*}%
Here we can argue that a bond market operates in the background and $%
r_{b}=r_{k}=r.$
\end{frame}

\begin{frame}
\frametitle{A storage economy with diminishing returns}
Consider the case above except that with the storage technology $k$ units
stored in period 1 can be transformed to $k^{\alpha }\;$units tomorrow with $%
0<\alpha <1$. Also, let $y_{2}=0$ and simplify notation so that $y_{1}=y.$ A
bond market still exists, but operates in the background. Let $k$ be the
amount of storage chosen by the representative agent. The first and second
period budget constraints are 
\begin{eqnarray*}
c_{1}+k &=&y \\
c_{2} &=&k^{\alpha }.
\end{eqnarray*}%
The individual's problem is to maximize 
\begin{equation*}
\ln \left( y-k\right) +\beta \ln \left( k^{\alpha }\right)
\end{equation*}%
so foc is $\frac{1}{y-k}=\frac{\beta \alpha k^{\alpha -1}}{k^{\alpha }}.$%
\end{frame}

\begin{frame}
\frametitle{Slide 79}
Then 
\begin{equation*}
\frac{1}{y-k}=\frac{\beta \alpha }{k}
\end{equation*}%
\begin{equation*}
k=\frac{\alpha \beta y}{\alpha \beta +1}.
\end{equation*}%
Show that 
\begin{eqnarray*}
R &=&\left( 1+r_{b}\right) =\alpha k^{\alpha -1} \\
r_{b} &=&\alpha \left( \frac{\alpha \beta y}{\alpha \beta +1}\right)
^{\alpha -1}-1.
\end{eqnarray*}%
$\allowbreak $

\begin{itemize}
\end{itemize}
\end{frame}

\begin{frame}
\frametitle{Note in all cases we have}
\begin{itemize}
\item Note in all cases we have 
\begin{equation*}
\frac{1}{p_{2}}=R=\frac{\frac{1}{c_{1}}}{\frac{\beta }{c_{2}}}=\frac{u_{1}}{%
\beta u_{2}}=\left\vert MRS_{1,2}\right\vert
\end{equation*}

\item In our first example, $p_{2},$ adjusted to make this hold.

\item In our second, the interest rate adjusted to make this hold.

\item In our third, consumption adjusted to make this hold.

\item In our fourth example, consumption and interest rates both adjust to
make this hold.
\end{itemize}
\end{frame}

\begin{frame}
\frametitle{A production economy with diminishing returns}
Consider the possibility of investment in physical capital. A representative
agent receives $y$ units of output in the first period of life and none in
the second period of life. The agent has one unit of time supplied
inelastically in the labor market in the second period of life. The agent
can save some consumption while young. This becomes capital which can be
rented to a firm in the second period. A representative firm can transform $%
k $ units of capital into $k^{\alpha }L^{1-\alpha }$ extra units of the good
in the next period. This firm hires capital at rate $r_{k}$ in the second
period and uses it to produce output. Capital is used up in the production
of the consumption good. That is, the depreciation rate is 1. In the second
period, the firm chooses $k$ to maximize profits%
\begin{equation*}
\pi =k^{\alpha }L^{1-\alpha }-r_{k}k-\omega L.
\end{equation*}%
Assume that the agent owns both the capital stock and the firm. The focs are:%
\begin{eqnarray*}
\alpha k^{\alpha -1}L^{1-\alpha } &=&r_{k} \\
\left( 1-\alpha \right) k^{\alpha }L^{-\alpha } &=&\omega .
\end{eqnarray*}%
\end{frame}

\begin{frame}
\frametitle{The representative agent chooses, and (savings) to}
The representative agent chooses $c_{1}$, $c_{2}$ and $s$ (savings) to
maximize 
\begin{equation*}
\underset{s,c_{1},c_{2}}{\max }\ln c_{1}+\beta \ln c_{2}
\end{equation*}%
subject to 
\begin{eqnarray*}
c_{1} &\leq &y-s \\
c_{2} &\leq &sR+\omega
\end{eqnarray*}%
where $R$ is the return to savings.
\end{frame}

\begin{frame}
\frametitle{Rewrite as}
Rewrite as%
\begin{equation*}
\underset{s}{\max }\ln \left( y-s\right) +\beta \ln \left( sR+\omega \right) 
\end{equation*}%
so that focs is%
\begin{equation*}
\frac{1}{y-s}=\frac{\beta R}{sR+\omega }.
\end{equation*}%
In equilibrium $R=r_{k}=\alpha k^{\alpha -1}L^{1-\alpha },$ $k=s,L=1$ 
\begin{equation*}
\frac{1}{y-k}=\frac{\beta \alpha k^{\alpha -1}L^{1-\alpha }}{k\alpha
k^{\alpha -1}L^{1-\alpha }+\left( 1-\alpha \right) k^{\alpha }L^{-\alpha }}
\end{equation*}%
or 
\begin{eqnarray*}
\frac{1}{y-k} &=&\frac{\beta \alpha }{k\alpha +\left( 1-\alpha \right) k} \\
k &=&\frac{\alpha \beta y}{\alpha \beta +1}.
\end{eqnarray*}%
and 
\begin{equation*}
c=y-\frac{\alpha \beta y}{\alpha \beta +1}=\frac{y}{\alpha \beta +1}
\end{equation*}
$\allowbreak $
\end{frame}

\begin{frame}
\frametitle{On your own...Because it is simpler and the proper conditions hold,...}
On your own...Because it is simpler and the proper conditions hold, we can
instead solve the social planner's problem. This is 
\begin{equation*}
\underset{k}{\max }\ln \left( c_{1}\right) +\beta \ln \left( c_{2}\right)
\end{equation*}%
subject to%
\begin{equation*}
c_{1}+k\leq y
\end{equation*}%
\begin{equation*}
c_{2}\leq k^{\alpha }
\end{equation*}%
since $L=1.$ We plug the feasibility constraint into the utility function to
get: 
\begin{equation*}
\underset{k}{\max }\ln \left( y-k\right) +\beta \ln \left( k^{\alpha
}\right) .
\end{equation*}
\end{frame}

\begin{frame}
\frametitle{The foc is}
The foc is 
\begin{equation*}
\frac{1}{y-k}=\beta \frac{\alpha k^{\alpha -1}}{k^{\alpha }}
\end{equation*}%
or%
\begin{equation*}
k=\frac{\alpha \beta y}{1+\alpha \beta }.
\end{equation*}%
i.e. a constant share is saved (since we have log utility). Now consumption
in each period is 
\begin{eqnarray*}
c_{1} &=&y-k=\frac{1}{1+\alpha \beta }y \\
c_{2} &=&\left( \frac{\alpha \beta y}{1+\alpha \beta }\right) ^{\alpha }.
\end{eqnarray*}
\end{frame}

\begin{frame}
\frametitle{An economy with three periods}
Consider the same problem except that the representative agent lives 3
periods. We solve the social planner's problem%
\begin{equation*}
\underset{c_{1},c_{2},c_{3},k_{2},k_{3}}{\max }\ln c_{1}+\beta \ln
c_{2}+\beta ^{2}\ln c_{3}
\end{equation*}%
subject to%
\begin{equation*}
c_{1}+k_{2}\leq y
\end{equation*}%
\begin{equation*}
c_{2}+k_{3}\leq k_{2}^{\alpha }L^{1-\alpha }
\end{equation*}%
\begin{equation*}
c_{3}\leq k_{3}^{\alpha }L^{1-\alpha }
\end{equation*}%
where $c_{3}$ is period 3 consumption, $k_{i}$ is capital available in
period $i$.
\end{frame}

\begin{frame}
\frametitle{All resource constraints will bind. Write as}
All resource constraints will bind. Write as 
\begin{equation*}
L=\ln \left( y-k_{2}\right) +\beta \ln \left( k_{2}^{\alpha }-k_{3}\right)
+\beta ^{2}\ln \left( k_{3}^{\alpha }\right)
\end{equation*}%
so that the first order conditions are%
\begin{eqnarray*}
\frac{1}{y-k_{2}} &=&\frac{\alpha \beta k_{2}^{\alpha -1}}{k_{2}^{\alpha
}-k_{3}} \\
\frac{\beta }{k_{2}^{\alpha }-k_{3}} &=&\frac{\beta ^{2}\alpha k_{3}^{\alpha
-1}}{k_{3}^{\alpha }}.
\end{eqnarray*}%
\end{frame}

\begin{frame}
\frametitle{From second}
From second 
\begin{equation*}
k_{3}=\alpha \beta \frac{k_{2}^{\alpha }}{\alpha \beta +1}
\end{equation*}%
and this into the first gives 
\begin{equation*}
\frac{1}{y-k_{2}}=\frac{\alpha \beta k_{2}^{\alpha -1}}{k_{2}^{\alpha
}-\alpha \beta \frac{k_{2}^{\alpha }}{\alpha \beta +1}}
\end{equation*}%
\begin{equation*}
k_{2}=\frac{y\left( \alpha \beta +\alpha ^{2}\beta ^{2}\right) }{1+\alpha
\beta +\alpha ^{2}\beta ^{2}}.
\end{equation*}%
Since $c_{1}=y-k_{2}$ we have%
\begin{equation*}
c_{1}=y-\frac{y\left( \alpha \beta +\alpha ^{2}\beta ^{2}\right) }{1+\alpha
\beta +\alpha ^{2}\beta ^{2}}
\end{equation*}%
so that 
\begin{equation*}
c_{1}=\frac{y}{1+\alpha \beta +\alpha ^{2}\beta ^{2}}.
\end{equation*}%
\end{frame}

\begin{frame}
\frametitle{Compare this to the two period case}
Compare this to the two period case%
\begin{equation*}
c_{1}=\frac{1}{1+\alpha \beta }y
\end{equation*}%
\begin{equation*}
c_{1}=\frac{y}{1+\alpha \beta +\alpha ^{2}\beta ^{2}}.
\end{equation*}%
How do you expect your answer to change if agents live four periods?
How do you expect your answer to change if agents live forever?
\end{frame}

\begin{frame}
\frametitle{An economy with heterogeneous agents}
Consider an economy with 2 agents. Agent 1 is endowed with $\alpha _{1}$
units of consumption in period 1 and $\theta _{1}$ units of consumption good
in period 2. Endowments for agent 2 are $\alpha _{2}$ and $\theta _{2}.$
Each agent has preferences given by 
\begin{equation*}
U_{i}\left( c_{i}\left( 1\right) ,c_{i}\left( 2\right) \right) =\ln
c_{i}\left( 1\right) +\beta \ln \left( c_{i}\left( 2\right) \right)
\end{equation*}%
where $c_{i}\left( j\right) $ is consumption in period $j$ by agent $i$ for $%
j,i\in \left\{ 1,2\right\} .$ Note $U_{i}$ is utility of agent $i$ (not a
derivative). Agents can issue or purchase 1 period bonds in the first period
that yield $1+r=R$ units of the consumption good in period 2 per unit of
bond purchased in period 1. There is no default on bonds.

Find consumption in period 1, consumption in period 2 and bond holdings in
period 1 for each agent as a function of parameters and $r.$\newline
\end{frame}

\begin{frame}
\frametitle{The problem for agent can be written as}
The problem for agent $i\in \left\{ 1,2\right\} $ can be written as%
\begin{eqnarray*}
&&\underset{c_{i}\left( 1\right) ,c_{i}\left( 2\right) ,b_{i}}{\max }\ln
c_{i}\left( 1\right) +\beta \ln \left( c_{i}\left( 2\right) \right) \\
&&\text{subject to} \\
c_{i}\left( 1\right) +b_{i} &=&\alpha _{i} \\
c_{i}\left( 2\right) &=&b_{i}R+\theta _{i}.
\end{eqnarray*}%
Solving gives 
\begin{equation*}
c_{i}\left( 1\right) =\frac{1}{1+\beta }\left( \alpha _{i}+\frac{\theta _{i}%
}{R}\right)
\end{equation*}%
\begin{equation*}
c_{i}\left( 2\right) =\frac{\beta R}{1+\beta }\left( \alpha _{i}+\frac{%
\theta _{i}}{R}\right)
\end{equation*}%
\begin{equation}
b_{i}=\alpha _{i}-\frac{1}{1+\beta }\left( \alpha _{i}+\frac{\theta _{i}}{R}%
\right) .  \label{bond}
\end{equation}
\end{frame}

\begin{frame}
\frametitle{Slide 92}
Now we find the maximum lifetime utility ($v_{ai}$) for each agent $i$ as a
function of $\alpha _{i}$, $\theta _{i}$ and $R$. \ Plugging $c_{i}\left(
1\right) $ and $c_{i}\left( 2\right) $ into the utility function gives%
\begin{equation*}
v_{ai}=\ln \frac{1}{1+\beta }\left( \alpha _{i}+\frac{\theta _{i}}{R}\right)
+\beta \ln \frac{\beta R}{1+\beta }\left( \alpha _{i}+\frac{\theta _{i}}{R}%
\right)
\end{equation*}%
or%
\begin{equation*}
v_{ai}=\ln \left[ \left( \frac{1}{1+\beta }\left( \alpha _{i}+\frac{\theta
_{i}}{R}\right) \right) ^{1+\beta }\left( \beta R\right) ^{^{\beta }}\right]
\end{equation*}%
How does $v_{ai}$ change with $R$ when $\alpha _{i}=0?$ when $\theta _{i}=0?$
Why?
\end{frame}

\begin{frame}
\frametitle{Define an equilibrium}
Define an equilibrium. \newline
(Roughly) An equilibrium is a price ($R$) and allocations $c_{1}\left(
1\right) ,$ $c_{1}\left( 2\right) $, $c_{2}\left( 1\right) ,$ $c_{2}\left(
2\right) $ such that each agent optimizes (giving equation (\ref{bond})),
the goods market clears in each period ($c_{1}\left( 1\right) +c_{2}\left(
1\right) =\alpha _{1}+\alpha _{2}$ and $c_{1}\left( 2\right) +c_{2}\left(
2\right) =\theta _{1}+\theta _{2}$) and the bond market clears in period 1 $%
\left( b_{1}+b_{2}=0\right) .$

Find the equilibrium interest rate. Bond market clearing requires $%
b_{1}+b_{2}=0$ or 
\begin{equation*}
\alpha _{1}+\alpha _{2}=\frac{1}{1+\beta }\left( \alpha _{1}+\frac{\theta
_{1}}{R}\right) +\frac{1}{1+\beta }\left( \alpha _{2}+\frac{\theta _{2}}{R}%
\right)
\end{equation*}%
\begin{equation*}
\beta \left( \alpha _{1}+\alpha _{2}\right) =\frac{\theta _{1}}{R}+\frac{%
\theta _{2}}{R}
\end{equation*}%
\begin{equation*}
R=\frac{\theta _{1}+\theta _{2}}{\beta \left( \alpha _{1}+\alpha _{2}\right) 
}.
\end{equation*}
\end{frame}

\begin{frame}
\frametitle{Agent 1 will save if bond holdings are positive, i.e. if}
Agent 1 will save if bond holdings are positive, i.e. if 
\begin{equation*}
b_{1}=\alpha _{1}-\frac{1}{1+\beta }\left( \alpha _{1}+\frac{\theta _{1}}{R}%
\right) >0
\end{equation*}%
or%
\begin{equation*}
\alpha _{1}>\frac{1}{1+\beta }\left( \alpha _{1}+\frac{\theta _{1}\beta
\left( \alpha _{1}+\alpha _{2}\right) }{\theta _{1}+\theta _{2}}\right)
\end{equation*}%
\begin{equation*}
\alpha _{1}>\frac{\theta _{1}\left( \alpha _{1}+\alpha _{2}\right) }{\theta
_{1}+\theta _{2}}
\end{equation*}%
\begin{equation*}
\frac{\alpha _{1}}{\theta _{1}}>\frac{\alpha _{2}}{\theta _{2}}.
\end{equation*}
\end{frame}

\begin{frame}
\frametitle{An economy with agents who differ in preferences}
Consider again an economy with 2 agents. Each agent is endowed with $\alpha $
units of consumption in period 1 and $\theta $ units of consumption good in
period 2. Agent 1 has preferences given by 
\begin{equation*}
U_{1}=\ln c_{1}\left( 1\right) +\beta \ln \left( c_{1}\left( 2\right) \right)
\end{equation*}%
and Agent 2 has preferences given by 
\begin{equation*}
U_{2}=c_{2}\left( 1\right) +\beta c_{2}\left( 2\right)
\end{equation*}%
where $c_{i}\left( j\right) $ is consumption in period $j$ by agent $i$ for $%
j,i\in \left\{ 1,2\right\} .$ Agents can issue or purchase 1 period bonds in
the first period that yield $R$ units of the consumption good in period 2
per unit of bond purchased in period 1. There is no default on bonds.

Find consumption in period 1, consumption in period 2 and bond holdings in
period 1 for each agent as a function of parameters and $R.$\newline
\end{frame}

\begin{frame}
\frametitle{The problem for agent is like that above problem so}
The problem for agent $1$ is like that above problem so 
\begin{equation}
b_{1}=\alpha -\frac{1}{1+\beta }\left( \alpha +\frac{\theta }{R}\right) .
\label{bond1}
\end{equation}%
Putting budget constraints ($c_{2}\left( 1\right) =\alpha -b_{2}$ and $%
c_{2}\left( 2\right) =b_{2}R+\theta $) into the utility function for second
agent gives 
\begin{equation*}
\alpha -b_{2}+\beta \left[ b_{2}R+\theta \right]
\end{equation*}%
or 
\begin{equation*}
\alpha +\beta \theta +\left( \beta R-1\right) b_{2}
\end{equation*}%
which is maximized at%
\begin{equation*}
b_{2}=\left\{ 
\begin{array}{ccc}
-\infty & \text{if} & \beta R<1 \\ 
\text{any quantity} & \text{if} & \beta R=1 \\ 
\infty & \text{if} & \beta R>1.%
\end{array}%
\right.
\end{equation*}
\end{frame}

\begin{frame}
\frametitle{Slide 97}
Define an equilibrium. (Roughly) An equilibrium is a price ($r$) and
allocations $c_{1}\left( 1\right) ,$ $c_{1}\left( 2\right) $, $c_{2}\left(
1\right) ,$ $c_{2}\left( 2\right) $ such that each agent optimizes, the
goods market clears in each period ($c_{1}\left( 1\right) +c_{2}\left(
1\right) =2\alpha $ and $c_{1}\left( 2\right) +c_{2}\left( 2\right) =2\theta 
$) and the bond market clears in period 1 $\left( b_{1}+b_{2}=0\right) .$
Find the equilibrium interest rate. Bond market clearing requires $%
b_{1}+b_{2}=0$ which can occur only if $\beta R=1.$ This into (\ref{bond1})
gives%
\begin{equation*}
b_{1}=\alpha -\frac{1}{1+\beta }\left( \alpha +\beta \theta \right)
\end{equation*}%
and by bond market clearing 
\begin{equation*}
b_{2}=\frac{1}{1+\beta }\left( \alpha +\beta \theta \right) -\alpha .
\end{equation*}%
Agent 2 will be a borrow if $b_{2}<0\ $or 
\begin{equation*}
\frac{1}{1+\beta }\left( \alpha +\beta \theta \right) <\alpha
\end{equation*}%
or%
\begin{equation*}
\theta <\alpha .
\end{equation*}%
If $\theta $ (second period endowment) is small, the agent with log
preferences wishes to save for purposes of consumption smoothing. \ The
other agent, indifferent to consumption smoothing, accommodates in
equilibrium and borrows.
\end{frame}

\begin{frame}
\frametitle{Overlapping generations}
\subsection{A first example}

The economy consists of overlapping generations of two-period-lived
homogeneous agents, a representative firm producing a single good using
physical capital and labor and a government.

\begin{itemize}
\item A representative agent is born each period. The population of each
generation is normalized to 1.

\item At each point in time there are young and old.

\item Discuss overlapping generations structure.

\item The initial old agents are endowed with $K_{0}$ units of capital$.$

\item All young have $\phi _{y}$ units of time and all old have $\phi _{o}$
units of time where $\phi _{y}+\phi _{o}=1.$ Time is supplied exogenously to
the labor market.
\end{itemize}
\end{frame}

\begin{frame}
\frametitle{Wage per unit of time is give by}
\begin{itemize}
\item Wage per unit of time is give by $\omega _{t}$.

\item Savings in period $t$ $(s_{t})$ returns $s_{t}R_{t+1}$ units of the
final good in period $t+1.$

\item The total capital stock evolves according to $K_{t+1}=K_{t}\left(
1-\delta \right) +I_{t}$ where $I_{t}$ is investment in capital by the young.

\item A unit of capital purchased in period $t$ returns $1-\delta +r_{t+1}$
units where $r_{t+1}$ is the period $t+1$ rental rate. We set $\delta =1.$

\item Government spends $\tilde{e}_{t}$ on wasteful expenditures.
\end{itemize}
\end{frame}

\begin{frame}
\frametitle{Definition. A competitive equilibrium in this environment is a}
\textbf{Definition}. A competitive equilibrium in this environment is a
sequence of consumption levels and portfolio holdings $%
\{c_{t,t},c_{t,t+1},s_{t}\}_{t=0}^{t=\infty }$ chosen by the representative
agent of each generation, a sequence of outputs and inputs chosen by the
representative firm in each period $\{Y_{t},K_{t},L_{t}\}_{t=0}^{t=\infty }$%
, a sequence of government policies $\tilde{\tau}_{y},\tilde{\tau}_{o},\{%
\tilde{e}_{t}\}_{t=0}^{t=\infty }$, a sequence of prices $\{\omega
_{t},r_{t}\}_{t=0}^{t=\infty }$, and an initial condition $\{K_{0}\}$ such
that the initial old consume $r_{0}K_{0}+\omega _{0}\phi _{o}-\tilde{\tau}%
_{o}$ and

\begin{enumerate}
\item a period $t$ young agents chooses $c_{t,t},c_{t,t+1}\ $and $s_{t}$ to
solve the agent's problem taking prices and government policy as given,

\item the firm chooses $Y_{t},K_{t}$ and $L_{t}$ in period $t$ to maximize
profits taking prices, government policy and the production possibilities as
given,

\item the government chooses $\{\tilde{\tau}_{y,t},\tilde{\tau}_{o,t},\tilde{%
e}_{t}\}$ subject to a balanced budget constraint.

\item the goods market clears: $Y_{t}=\tilde{e}_{t}+c_{t,t}+c_{t-1,t}+I_{t}.$

\item the capital market clears, $K_{t}=I_{t-1}=s_{t-1}.$

\item the labor market clears: $L_{t}=L=\phi _{y}+\phi _{o}=1.$

\item No arbitrage requires $r_{t+1}=R_{t+1}$.\newline
\end{enumerate}

(We have $s_{t}$ instead of $s_{t,t}$ since only the young save so the extra 
$t$ is not needed. If agents lived 3 periods and saved in the first and
second period of life we would need $s_{t,t}$ and $s_{t,t+1}.)$
\end{frame}

\begin{frame}
\frametitle{A steady state satisfies the definition above and has the following}
A \textbf{steady state} satisfies the definition above and has the following
additional properties:\newline
(i) government policy is time invariant, $\tilde{e}_{t}=\tilde{e}$ \newline
(ii) the stock of physical capital, labor supply, wages, interest rates,
consumption by young, consumption by the old, output, and government
expenditures are all constant.

That is

\begin{itemize}
\item $c_{t,t}=c_{t+1,t+1}=c_{t+2,t+2}...=c_{y}$

\item $c_{t,t+1}=c_{t+1,t++2}=c_{t+2,t+3}...=c_{0}$
\end{itemize}
\end{frame}

\begin{frame}
\frametitle{Government spends a share of output on wasteful expenditures so that}
Government spends a share $e_{t}$ of output on wasteful expenditures so that 
$\tilde{e}_{t}=e_{t}Y_{t}$. Since this is the only type of spending for
government, a balanced budget in the steady state requires 
\begin{equation*}
\tilde{\tau}_{y}+\tilde{\tau}_{o}=eY.
\end{equation*}%
Agents take taxes as given. However, we assume (since it makes our math
easier) that these are fixed shares of the wage per unit of time in
equilibrium so that $\tilde{\tau}_{y}=\tau _{y}\omega $ and $\tilde{\tau}%
_{o}=\tau _{o}\omega $. Given this, in a steady state we have 
\begin{equation*}
\omega \left( \tau _{y}+\tau _{o}\right) =eY.
\end{equation*}%
\end{frame}

\begin{frame}
\frametitle{The representative agent{'}s problem is}
The representative agent's problem is: 
\begin{equation*}
\underset{c_{t,t},c_{t,t+1},s_{t}}{\max }\ln c_{t,t}+\beta \ln c_{t,t+1}
\end{equation*}%
subject to 
\begin{eqnarray*}
c_{t,t}+s_{t}+\tilde{\tau}_{y} &\leq &\phi _{y}\omega _{t} \\
c_{t,t+1}+\tilde{\tau}_{o} &\leq &R_{t+1}s_{t}+\phi _{o}\omega _{t+1}.
\end{eqnarray*}%
Case 1: Assume no bond market or storage. $R_{t+1}=0.$ $Y_{t}=\mu _{\ell
}L_{t}$.

In this case $R_{t+1}=r_{t+1}=0$ and $\omega _{t}=\omega _{t+1}=\mu _{\ell
}. $ As no bond market can operate%
\begin{eqnarray*}
s_{t} &=&0 \\
c_{t,t} &=&\phi _{y}\omega _{t}-\tilde{\tau}_{y} \\
c_{t,t+1} &=&\phi _{o}\omega _{t+1}-\tilde{\tau}_{o}
\end{eqnarray*}%
\end{frame}

\begin{frame}
\frametitle{Slide 104}
Consider the steady state and suppose $e=0$ so $\tilde{\tau}_{y}=-\tilde{\tau%
}_{o}=\tilde{\tau}.$ Then steady state utility is%
\begin{equation*}
\ln \left( \phi _{y}\omega -\tilde{\tau}\right) +\beta \ln \left( \phi
_{o}\omega +\tilde{\tau}\right)
\end{equation*}
\begin{equation*}
\ln \left( \phi _{y}\omega -\tau \omega \right) +\beta \ln \left( \phi
_{o}\omega +\tau \omega \right)
\end{equation*}
\begin{equation*}
\ln \left( \phi _{y}\mu _{\ell }-\tau \mu _{\ell }\right) +\beta \ln \left(
\phi _{o}\mu _{\ell }+\tau \mu _{\ell }\right)
\end{equation*}%
Suppose government wants to set $\tau $ to maximize steady state utility.
i.e. 
\begin{equation*}
\underset{\tau }{\max }\ln \left( \phi _{y}\mu _{\ell }-\mu _{\ell }\tau
\right) +\beta \ln \left( \phi _{o}\mu _{\ell }+\mu _{\ell }\tau \right)
\end{equation*}%
Foc is 
\begin{equation*}
\frac{1}{\phi _{y}-\tau }=\frac{\beta }{\phi _{o}+\tau }
\end{equation*}%
\begin{equation*}
\tau =\frac{\beta \phi _{y}-\phi _{0}}{1+\beta }
\end{equation*}%
So $\tau >0$ iff%
\begin{equation*}
\phi _{y}>\frac{\phi _{o}}{\beta }
\end{equation*}
\end{frame}

\begin{frame}
\frametitle{Slide 105}
Case 2: $R_{t+1}>0,$ $Y_{t}=\mu _{\ell }L_{t}+\mu _{k}K_{t}$. In this case $%
R_{t+1}=r_{t+1}=\mu _{k}$ and $\omega _{t+1}=\omega _{t}=\mu _{\ell }.$ The
agent's problem is 
\begin{eqnarray*}
&&\underset{c_{t,t},c_{t,t+1},s_{t}}{\max }\ln c_{t,t}+\beta \ln c_{t,t+1} \\
&&+\lambda _{t}\left( \phi _{y}\omega _{t}-c_{t,t}-\tilde{\tau}%
_{y}-s_{t}\right) \\
&&+\lambda _{t+1}\left( R_{t+1}s_{t}+\phi \omega _{t+1}-c_{t,t+1}-\tilde{\tau%
}_{o}\right)
\end{eqnarray*}%
Optimal consumption requires 
\begin{eqnarray*}
\frac{1}{c_{t,t}} &=&\lambda _{t} \\
\frac{\beta }{c_{t,t+1}} &=&\lambda _{t+1} \\
\lambda _{t} &=&\lambda _{t+1}r_{t+1}
\end{eqnarray*}%
(since $r_{t}=R_{t}$).
\end{frame}

\begin{frame}
\frametitle{The first and second into the third gives}
The first and second into the third gives%
\begin{equation*}
c_{t,t+1}=\beta r_{t+1}c_{t,t}
\end{equation*}%
Since budget constraints hold with equality we can treat them as a single
lifetime budget constraint.%
\begin{equation*}
c_{t,t}+\frac{c_{t,t+1}}{r_{t+1}}+\tilde{\tau}_{y}+\frac{\tilde{\tau}_{o}}{%
r_{t+1}}=\omega _{t}\phi _{y}+\frac{\phi _{o}\omega _{t+1}}{r_{t+1}}
\end{equation*}%
plugging in for $c_{t,t+1}$ and considering the state we have: 
\begin{equation*}
c_{t,t}=c_{y}=\frac{1}{1+\beta }\left( \phi _{y}\omega -\tilde{\tau}_{y}+%
\frac{\phi _{o}\omega }{r}-\frac{\tilde{\tau}_{o}}{r}\right)
\end{equation*}%
\begin{equation*}
c_{t+1,t}=c_{o}=\frac{\beta R}{1+\beta }\left( \phi _{y}\omega -\tilde{\tau}%
_{y}+\frac{\phi _{o}\omega }{r}-\frac{\tilde{\tau}_{o}}{r}\right)
\end{equation*}%
and then since $s=\omega \phi _{y}-\tilde{\tau}_{y}-c_{t,t}$ 
\begin{equation*}
s_{t}=s=\phi _{y}\omega _{t}-\tilde{\tau}_{y}-\frac{1}{1+\beta }\left( \phi
_{y}\omega +\frac{\phi _{o}\omega }{r}-\tilde{\tau}_{y}-\frac{\tilde{\tau}%
_{o}}{r}\right)
\end{equation*}%
In a steady state and with again $\tilde{\tau}_{y}=-\tilde{\tau}_{o}=\tilde{%
\tau}$ 
\begin{equation*}
K=\phi _{y}\omega -\tau \omega -\frac{1}{1+\beta }\left( \phi _{y}\omega +%
\frac{\phi _{o}\omega }{r}-\tau \omega +\frac{\tau \omega }{r}\right)
\end{equation*}%
\begin{equation*}
K=\phi _{y}\mu _{\ell }-\tau \mu _{\ell }-\frac{1}{1+\beta }\left( \phi
_{y}\mu _{\ell }+\frac{\phi _{o}\mu _{\ell }}{\mu _{k}}-\tau \mu _{\ell }+%
\frac{\tau \mu _{\ell }}{\mu _{k}}\right)
\end{equation*}%
\end{frame}

\begin{frame}
\frametitle{Equilibrium utility is}
Equilibrium utility is%
\begin{equation*}
\ln c_{y}+\beta \ln c_{o}
\end{equation*}%
which we found to be the same as%
\begin{equation*}
\ln c_{y}+\beta \ln \beta rc_{y}
\end{equation*}%
Write this as 
\begin{equation*}
\left( 1+\beta \right) \ln c_{y}+\beta \ln \beta r.
\end{equation*}%
With $R=r$ unaffected by the level of capital, and hence saving, and hence
consumption, the level of $\tau $ that maximizes this is the same as the
level that maximizes 
\begin{equation*}
c_{y}=\frac{1}{1+\beta }\left( \phi _{y}\mu _{\ell }+\frac{\phi _{o}\mu
_{\ell }}{\mu _{k}}-\tau \mu _{\ell }+\frac{\tau \mu _{\ell }}{\mu _{k}}%
\right)
\end{equation*}%
The government should set $\tau >0$ iff%
\begin{equation*}
-\mu _{\ell }+\frac{\mu _{\ell }}{\mu _{k}}>0
\end{equation*}%
\begin{equation*}
\mu _{k}<1
\end{equation*}

Intuition?
\end{frame}

\begin{frame}
\frametitle{Case 3: . In this case, in equilibrium}
Case 3: $R_{t+1}>0,$ $Y_{t}=K_{t}^{\alpha }L_{t}^{1-\alpha }$.

In this case, in equilibrium%
\begin{eqnarray*}
r_{t} &=&\alpha K_{t}^{\alpha -1} \\
\omega _{t} &=&\left( 1-\alpha \right) K_{t}^{\alpha }.
\end{eqnarray*}%
As before we have 
\begin{equation*}
s=\phi _{y}\omega _{t}-\tilde{\tau}_{y}-\frac{1}{1+\beta }\left( \phi
_{y}\omega _{t}+\frac{\phi _{o}\omega _{t+1}}{r_{t+1}}-\tilde{\tau}_{y}-%
\frac{\tilde{\tau}_{o}}{r_{t+1}}\right)
\end{equation*}%
or in a steady state equilibrium 
\begin{equation*}
K=\phi _{y}\omega -\tau \omega -\frac{1}{1+\beta }\left( \phi _{y}\omega +%
\frac{\phi _{o}\omega }{r}-\tau \omega +\frac{\tau \omega }{r}\right)
\end{equation*}%
Lets write this as 
\begin{equation*}
\frac{K}{\omega }=\phi _{y}-\tau -\frac{1}{1+\beta }\left( \phi _{y}-\tau
+\left( \frac{\phi _{o}+\tau }{r}\right) \right)
\end{equation*}%
\end{frame}

\begin{frame}
\frametitle{Slide 109}
Note 
\begin{equation*}
\frac{K}{\omega }=\frac{\alpha }{r\left( 1-\alpha \right) }
\end{equation*}%
so this is 
\begin{equation*}
\frac{\alpha }{r\left( 1-\alpha \right) }=\left( \phi _{y}-\tau \right) -%
\frac{1}{1+\beta }\left( \left( \phi _{y}-\tau \right) +\left( \frac{\phi
_{o}+\tau }{r}\right) \right)
\end{equation*}%
or%
\begin{equation*}
\frac{1}{r}\left( \frac{\alpha }{1-\alpha }+\frac{\phi _{o}+\tau }{1+\beta }%
\right) =\frac{\beta \left( \phi _{y}-\tau \right) }{1+\beta }
\end{equation*}%
or%
\begin{equation*}
\frac{1}{\alpha K^{\alpha -1}}=\frac{\beta \left( \phi _{y}-\tau \right) }{%
\left( 1+\beta \right) \left( \frac{\alpha }{1-\alpha }+\frac{\phi _{o}+\tau 
}{1+\beta }\right) }
\end{equation*}%
giving 
\begin{equation*}
K=\left( \frac{\beta \alpha \left( \phi _{y}-\tau \right) }{\left( 1+\beta
\right) \left( \frac{\alpha }{1-\alpha }+\frac{\phi _{o}+\tau }{1+\beta }%
\right) }\right) ^{\frac{1}{1-\alpha }}.
\end{equation*}
\end{frame}

\begin{frame}
\frametitle{We can see that taxation decreases savings. The effect of taxation on}
We can see that taxation decreases savings. The effect of taxation on
utility is tough. We know from before that utility is 
\begin{equation*}
\left( 1+\beta \right) \ln c_{y}+\beta \ln \beta r.
\end{equation*}%
But (after cleaning things up a bit)%
\begin{equation*}
c_{y}=\frac{1}{1+\beta }\left( \left( 1-\alpha \right) K^{\alpha }\left(
\phi _{y}-\tau \right) +\frac{\left( \phi _{o}+\tau \right) \left( 1-\alpha
\right) K}{\alpha }\right)
\end{equation*}%
so 
\begin{equation*}
\frac{\partial \ln c_{y}}{\partial K}\frac{\partial K}{\partial \tau }
\end{equation*}%
is already tough and then we need to also consider 
\begin{equation*}
\frac{\partial \ln r}{\partial K}\frac{\partial K}{\partial \tau }.
\end{equation*}%
As such, we will not consider what tax rate maximizes utility as in the
earlier cases.
\end{frame}

\end{document}
